\begingroup\tcbset{breakable/.default=false}
\renewcommand{\Needspace}[1]{\par\begingroup
\dimen0=#1\relax\vskip0pt plus\dimen0\penalty-100
\vskip0pt plus-\dimen0\vskip\dimen0\penalty9999\vskip-\dimen0
\vskip0pt\endgroup}
\chapter{几何光学：光线、反射和折射}


\graphicspath{{/Users/dongxingxi/Documents/物理竞赛一轮资料/sample/xes-original/forty-second-section/}{/Users/dongxingxi/Documents/物理竞赛一轮资料/sample/xes-original/forty-second-section/assets/}}


\section{光线模型、波面与光学发展}


\newcommand{\AtmosphereWavefrontPanel}{%
\begin{center}\begin{minipage}{\linewidth}\centering
\begin{tikzpicture}[>=Latex,font=\small,every node/.style={fill=none}]
\draw[thick] plot coordinates {(2.8670322,-1.4924825) (2.8279646,-1.4378127) (2.7878547,-1.3839031) (2.7467173,-1.3307734) (2.7045675,-1.2784432) (2.6614209,-1.2269318) (2.6172933,-1.1762583) (2.5722011,-1.1264412) (2.5261608,-1.0774989) (2.4791894,-1.0294495) (2.4313043,-0.9823107) (2.3825230,-0.9360999) (2.3328636,-0.8908340) (2.2823444,-0.8465298) (2.2309840,-0.8032037) (2.1788012,-0.7608715) (2.1258154,-0.7195488) (2.0720461,-0.6792510) (2.0175131,-0.6399929) (1.9622365,-0.6017888) (1.9062366,-0.5646530) (1.8495342,-0.5285991) (1.7921500,-0.4936404) (1.7341053,-0.4597897) (1.6754215,-0.4270596) (1.6161201,-0.3954621) (1.5562231,-0.3650088) (1.4957525,-0.3357111) (1.4347306,-0.3075796) (1.3731799,-0.2806248) (1.3111231,-0.2548565) (1.2485830,-0.2302843) (1.1855827,-0.2069173) (1.1221455,-0.1847640) (1.0582946,-0.1638327) (0.9940537,-0.1441309) (0.9294464,-0.1256661) (0.8644965,-0.1084449) (0.7992280,-0.0924738) (0.7336650,-0.0777587) (0.6678315,-0.0643049) (0.6017519,-0.0521174) (0.5354504,-0.0412007) (0.4689516,-0.0315588) (0.4022800,-0.0231953) (0.3354601,-0.0161133) (0.2685166,-0.0103154) (0.2014741,-0.0058036) (0.1343573,-0.0025798) (0.0671910,-0.0006450) (0.0000000,0.0000000) (-0.0671910,-0.0006450) (-0.1343573,-0.0025798) (-0.2014741,-0.0058036) (-0.2685166,-0.0103154) (-0.3354601,-0.0161133) (-0.4022800,-0.0231953) (-0.4689516,-0.0315588) (-0.5354504,-0.0412007) (-0.6017519,-0.0521174) (-0.6678315,-0.0643049) (-0.7336650,-0.0777587) (-0.7992280,-0.0924738) (-0.8644965,-0.1084449) (-0.9294464,-0.1256661) (-0.9940537,-0.1441309) (-1.0582946,-0.1638327) (-1.1221455,-0.1847640) (-1.1855827,-0.2069173) (-1.2485830,-0.2302843) (-1.3111231,-0.2548565) (-1.3731799,-0.2806248) (-1.4347306,-0.3075796) (-1.4957525,-0.3357111) (-1.5562231,-0.3650088) (-1.6161201,-0.3954621) (-1.6754215,-0.4270596) (-1.7341053,-0.4597897) (-1.7921500,-0.4936404) (-1.8495342,-0.5285991) (-1.9062366,-0.5646530) (-1.9622365,-0.6017888) (-2.0175131,-0.6399929) (-2.0720461,-0.6792510) (-2.1258154,-0.7195488) (-2.1788012,-0.7608715) (-2.2309840,-0.8032037) (-2.2823444,-0.8465298) (-2.3328636,-0.8908340) (-2.3825230,-0.9360999) (-2.4313043,-0.9823107) (-2.4791894,-1.0294495) (-2.5261608,-1.0774989) (-2.5722011,-1.1264412) (-2.6172933,-1.1762583) (-2.6614209,-1.2269318) (-2.7045675,-1.2784432) (-2.7467173,-1.3307734) (-2.7878547,-1.3839031) (-2.8279646,-1.4378127) (-2.8670322,-1.4924825)};
\draw[thick] plot coordinates {(3.1127778,-1.3204095) (3.0703616,-1.2610538) (3.0268137,-1.2025234) (2.9821502,-1.1448397) (2.9363876,-1.0880240) (2.8895427,-1.0320974) (2.8416327,-0.9770804) (2.7926755,-0.9229933) (2.7426889,-0.8698559) (2.6916914,-0.8176880) (2.6397018,-0.7665088) (2.5867393,-0.7163370) (2.5328234,-0.6671912) (2.4779739,-0.6190895) (2.4222112,-0.5720497) (2.3655556,-0.5260890) (2.3080282,-0.4812244) (2.2496501,-0.4374725) (2.1904428,-0.3948494) (2.1304282,-0.3533707) (2.0696283,-0.3130518) (2.0080657,-0.2739076) (1.9457629,-0.2359524) (1.8827429,-0.1992002) (1.8190291,-0.1636647) (1.7546447,-0.1293588) (1.6896137,-0.0962953) (1.6239599,-0.0644863) (1.5577075,-0.0339436) (1.4908810,-0.0046783) (1.4235051,0.0232986) (1.3556044,0.0499770) (1.2872041,0.0753469) (1.2183294,0.0993991) (1.1490056,0.1221245) (1.0792583,0.1435150) (1.0091132,0.1635625) (0.9385962,0.1822598) (0.8677333,0.1995998) (0.7965505,0.2155763) (0.7250742,0.2301833) (0.6533306,0.2434154) (0.5813462,0.2552678) (0.5091475,0.2657362) (0.4367612,0.2748165) (0.3642139,0.2825056) (0.2915323,0.2888005) (0.2187433,0.2936989) (0.1458737,0.2971991) (0.0729503,0.2992997) (0.0000000,0.3000000) (-0.0729503,0.2992997) (-0.1458737,0.2971991) (-0.2187433,0.2936989) (-0.2915323,0.2888005) (-0.3642139,0.2825056) (-0.4367612,0.2748165) (-0.5091475,0.2657362) (-0.5813462,0.2552678) (-0.6533306,0.2434154) (-0.7250742,0.2301833) (-0.7965505,0.2155763) (-0.8677333,0.1995998) (-0.9385962,0.1822598) (-1.0091132,0.1635625) (-1.0792583,0.1435150) (-1.1490056,0.1221245) (-1.2183294,0.0993991) (-1.2872041,0.0753469) (-1.3556044,0.0499770) (-1.4235051,0.0232986) (-1.4908810,-0.0046783) (-1.5577075,-0.0339436) (-1.6239599,-0.0644863) (-1.6896137,-0.0962953) (-1.7546447,-0.1293588) (-1.8190291,-0.1636647) (-1.8827429,-0.1992002) (-1.9457629,-0.2359524) (-2.0080657,-0.2739076) (-2.0696283,-0.3130518) (-2.1304282,-0.3533707) (-2.1904428,-0.3948494) (-2.2496501,-0.4374725) (-2.3080282,-0.4812244) (-2.3655556,-0.5260890) (-2.4222112,-0.5720497) (-2.4779739,-0.6190895) (-2.5328234,-0.6671912) (-2.5867393,-0.7163370) (-2.6397018,-0.7665088) (-2.6916914,-0.8176880) (-2.7426889,-0.8698559) (-2.7926755,-0.9229933) (-2.8416327,-0.9770804) (-2.8895427,-1.0320974) (-2.9363876,-1.0880240) (-2.9821502,-1.1448397) (-3.0268137,-1.2025234) (-3.0703616,-1.2610538) (-3.1127778,-1.3204095)};
\draw[thick] plot coordinates {(3.3585234,-1.1483366) (3.3127585,-1.0842949) (3.2657727,-1.0211436) (3.2175831,-0.9589060) (3.1682077,-0.8976049) (3.1176645,-0.8372630) (3.0659722,-0.7779025) (3.0131498,-0.7195454) (2.9592169,-0.6622130) (2.9041933,-0.6059266) (2.8480993,-0.5507068) (2.7909556,-0.4965741) (2.7327831,-0.4435484) (2.6736035,-0.3916492) (2.6134384,-0.3408957) (2.5523100,-0.2913066) (2.4902409,-0.2429001) (2.4272540,-0.1956940) (2.3633725,-0.1497059) (2.2986199,-0.1049526) (2.2330200,-0.0614507) (2.1665972,-0.0192161) (2.0993757,0.0217356) (2.0313805,0.0613892) (1.9626366,0.0997302) (1.8931693,0.1367444) (1.8230042,0.1724182) (1.7521672,0.2067385) (1.6806844,0.2396925) (1.6085822,0.2712681) (1.5358870,0.3014538) (1.4626258,0.3302383) (1.3888255,0.3576112) (1.3145133,0.3835621) (1.2397166,0.4080817) (1.1644629,0.4311609) (1.0887801,0.4527912) (1.0126960,0.4729645) (0.9362386,0.4916735) (0.8594361,0.5089113) (0.7823169,0.5246715) (0.7049093,0.5389482) (0.6272419,0.5517364) (0.5493434,0.5630311) (0.4712423,0.5728283) (0.3929676,0.5811244) (0.3145480,0.5879163) (0.2360125,0.5932014) (0.1573900,0.5969780) (0.0787095,0.5992444) (0.0000000,0.6000000) (-0.0787095,0.5992444) (-0.1573900,0.5969780) (-0.2360125,0.5932014) (-0.3145480,0.5879163) (-0.3929676,0.5811244) (-0.4712423,0.5728283) (-0.5493434,0.5630311) (-0.6272419,0.5517364) (-0.7049093,0.5389482) (-0.7823169,0.5246715) (-0.8594361,0.5089113) (-0.9362386,0.4916735) (-1.0126960,0.4729645) (-1.0887801,0.4527912) (-1.1644629,0.4311609) (-1.2397166,0.4080817) (-1.3145133,0.3835621) (-1.3888255,0.3576112) (-1.4626258,0.3302383) (-1.5358870,0.3014538) (-1.6085822,0.2712681) (-1.6806844,0.2396925) (-1.7521672,0.2067385) (-1.8230042,0.1724182) (-1.8931693,0.1367444) (-1.9626366,0.0997302) (-2.0313805,0.0613892) (-2.0993757,0.0217356) (-2.1665972,-0.0192161) (-2.2330200,-0.0614507) (-2.2986199,-0.1049526) (-2.3633725,-0.1497059) (-2.4272540,-0.1956940) (-2.4902409,-0.2429001) (-2.5523100,-0.2913066) (-2.6134384,-0.3408957) (-2.6736035,-0.3916492) (-2.7327831,-0.4435484) (-2.7909556,-0.4965741) (-2.8480993,-0.5507068) (-2.9041933,-0.6059266) (-2.9592169,-0.6622130) (-3.0131498,-0.7195454) (-3.0659722,-0.7779025) (-3.1176645,-0.8372630) (-3.1682077,-0.8976049) (-3.2175831,-0.9589060) (-3.2657727,-1.0211436) (-3.3127585,-1.0842949) (-3.3585234,-1.1483366)};
\fill (0.0000000,0.0000000) circle (1.2pt);
\node[below left,fill=none,inner sep=2pt] at (0.0000000,0.0000000) {$O$};
\draw[gray,dashed,thin] (-3.0000000,0.0000000)--(3.0000000,0.0000000);
\node[right,fill=none,inner sep=2pt] at (3.2000000,-0.2000000) {局地水平};
\draw[DeepBlue,thick,<-] plot coordinates {(0.0000000,0.0000000) (0.0350000,0.0155563) (0.0700000,0.0307340) (0.1050000,0.0455462) (0.1400000,0.0600054) (0.1750000,0.0741238) (0.2100000,0.0879130) (0.2450000,0.1013845) (0.2800000,0.1145492) (0.3150000,0.1274176) (0.3500000,0.1399998) (0.3850000,0.1523058) (0.4200000,0.1643450) (0.4550000,0.1761267) (0.4900000,0.1876596) (0.5250000,0.1989523) (0.5600000,0.2100131) (0.5950000,0.2208500) (0.6300000,0.2314706) (0.6650000,0.2418825) (0.7000000,0.2520927) (0.7350000,0.2621083) (0.7700000,0.2719358) (0.8050000,0.2815819) (0.8400000,0.2910526) (0.8750000,0.3003542) (0.9100000,0.3094922) (0.9450000,0.3184725) (0.9800000,0.3273004) (1.0150000,0.3359812) (1.0500000,0.3445199) (1.0850000,0.3529214) (1.1200000,0.3611905) (1.1550000,0.3693316) (1.1900000,0.3773492) (1.2250000,0.3852475) (1.2600000,0.3930307) (1.2950000,0.4007027) (1.3300000,0.4082673) (1.3650000,0.4157282) (1.4000000,0.4230890) (1.4350000,0.4303531) (1.4700000,0.4375238) (1.5050000,0.4446045) (1.5400000,0.4515980) (1.5750000,0.4585076) (1.6100000,0.4653360) (1.6450000,0.4720860) (1.6800000,0.4787603) (1.7150000,0.4853616) (1.7500000,0.4918923) (1.7850000,0.4983549) (1.8200000,0.5047518) (1.8550000,0.5110851) (1.8900000,0.5173570) (1.9250000,0.5235698) (1.9600000,0.5297253) (1.9950000,0.5358256) (2.0300000,0.5418726) (2.0650000,0.5478681) (2.1000000,0.5538139) (2.1350000,0.5597117) (2.1700000,0.5655632) (2.2050000,0.5713698) (2.2400000,0.5771333) (2.2750000,0.5828550) (2.3100000,0.5885364) (2.3450000,0.5941789) (2.3800000,0.5997838) (2.4150000,0.6053525) (2.4500000,0.6108861) (2.4850000,0.6163858) (2.5200000,0.6218529) (2.5550000,0.6272885) (2.5900000,0.6326936) (2.6250000,0.6380693) (2.6600000,0.6434166) (2.6950000,0.6487364) (2.7300000,0.6540298) (2.7650000,0.6592976) (2.8000000,0.6645408) (2.8350000,0.6697601) (2.8700000,0.6749563) (2.9050000,0.6801304) (2.9400000,0.6852830) (2.9750000,0.6904148) (3.0100000,0.6955267) (3.0450000,0.7006192) (3.0800000,0.7056930) (3.1150000,0.7107489) (3.1500000,0.7157873) (3.1850000,0.7208090) (3.2200000,0.7258144) (3.2550000,0.7308041) (3.2900000,0.7357788) (3.3250000,0.7407388) (3.3600000,0.7456847) (3.3950000,0.7506170) (3.4300000,0.7555362) (3.4650000,0.7604426) (3.5000000,0.7653368)};
\draw[gray,dashed,thin] (0.0000000,0.0000000)--(3.5000000,1.5750000);
\draw[gray,dashed,thin] (0.0000000,0.0000000)--(3.5000000,0.4550000);
\node[right,fill=none,inner sep=2pt] at (3.5000000,1.5800000) {视方向};
\node[right,fill=none,inner sep=2pt] at (3.5000000,0.3800000) {真实方向};
\node[below,fill=none,inner sep=2pt] at (-1.5000000,-0.4000000) {地表};
\begin{scope}[shift={(-4.0000000,2.5000000)}]
\draw[thick] plot coordinates {(0.2892544,-0.3447200) (0.2952266,-0.3396193) (0.3011088,-0.3344152) (0.3068993,-0.3291092) (0.3125963,-0.3237029) (0.3181981,-0.3181981) (0.3237029,-0.3125963) (0.3291092,-0.3068993) (0.3344152,-0.3011088) (0.3396193,-0.2952266) (0.3447200,-0.2892544) (0.3497157,-0.2831942) (0.3546048,-0.2770477) (0.3593860,-0.2708168) (0.3640576,-0.2645034) (0.3686184,-0.2581094) (0.3730669,-0.2516368) (0.3774018,-0.2450876) (0.3816216,-0.2384637) (0.3857253,-0.2317671) (0.3897114,-0.2250000) (0.3935789,-0.2181643) (0.3973264,-0.2112622) (0.4009529,-0.2042957) (0.4044573,-0.1972670) (0.4078385,-0.1901782) (0.4110955,-0.1830315) (0.4142272,-0.1758290) (0.4172327,-0.1685730) (0.4201112,-0.1612656) (0.4228617,-0.1539091) (0.4254834,-0.1465057) (0.4279754,-0.1390576) (0.4303371,-0.1315673) (0.4325678,-0.1240368) (0.4346666,-0.1164686) (0.4366331,-0.1088649) (0.4384665,-0.1012280) (0.4401664,-0.0935603) (0.4417322,-0.0858640) (0.4431635,-0.0781417) (0.4444598,-0.0703955) (0.4456206,-0.0626279) (0.4466458,-0.0548412) (0.4475349,-0.0470378) (0.4482876,-0.0392201) (0.4489038,-0.0313904) (0.4493833,-0.0235512) (0.4497259,-0.0157048) (0.4499315,-0.0078536) (0.4500000,0.0000000) (0.4499315,0.0078536) (0.4497259,0.0157048) (0.4493833,0.0235512) (0.4489038,0.0313904) (0.4482876,0.0392201) (0.4475349,0.0470378) (0.4466458,0.0548412) (0.4456206,0.0626279) (0.4444598,0.0703955) (0.4431635,0.0781417) (0.4417322,0.0858640) (0.4401664,0.0935603) (0.4384665,0.1012280) (0.4366331,0.1088649) (0.4346666,0.1164686) (0.4325678,0.1240368) (0.4303371,0.1315673) (0.4279754,0.1390576) (0.4254834,0.1465057) (0.4228617,0.1539091) (0.4201112,0.1612656) (0.4172327,0.1685730) (0.4142272,0.1758290) (0.4110955,0.1830315) (0.4078385,0.1901782) (0.4044573,0.1972670) (0.4009529,0.2042957) (0.3973264,0.2112622) (0.3935789,0.2181643) (0.3897114,0.2250000) (0.3857253,0.2317671) (0.3816216,0.2384637) (0.3774018,0.2450876) (0.3730669,0.2516368) (0.3686184,0.2581094) (0.3640576,0.2645034) (0.3593860,0.2708168) (0.3546048,0.2770477) (0.3497157,0.2831942) (0.3447200,0.2892544) (0.3396193,0.2952266) (0.3344152,0.3011088) (0.3291092,0.3068993) (0.3237029,0.3125963) (0.3181981,0.3181981) (0.3125963,0.3237029) (0.3068993,0.3291092) (0.3011088,0.3344152) (0.2952266,0.3396193) (0.2892544,0.3447200)};
\draw[thick] plot coordinates {(0.5463695,-0.6511378) (0.5576502,-0.6415031) (0.5687610,-0.6316731) (0.5796986,-0.6216506) (0.5904596,-0.6114388) (0.6010408,-0.6010408) (0.6114388,-0.5904596) (0.6216506,-0.5796986) (0.6316731,-0.5687610) (0.6415031,-0.5576502) (0.6511378,-0.5463695) (0.6605741,-0.5349223) (0.6698091,-0.5233123) (0.6788402,-0.5115428) (0.6876644,-0.4996175) (0.6962792,-0.4875400) (0.7046819,-0.4753140) (0.7128700,-0.4629432) (0.7208409,-0.4504314) (0.7285922,-0.4377824) (0.7361216,-0.4250000) (0.7434268,-0.4120882) (0.7505055,-0.3990508) (0.7573555,-0.3858919) (0.7639749,-0.3726155) (0.7703616,-0.3592255) (0.7765136,-0.3457261) (0.7824291,-0.3321215) (0.7881063,-0.3184156) (0.7935434,-0.3046128) (0.7987387,-0.2907171) (0.8036908,-0.2767329) (0.8083980,-0.2626644) (0.8128590,-0.2485159) (0.8170724,-0.2342918) (0.8210370,-0.2199962) (0.8247514,-0.2056336) (0.8282146,-0.1912084) (0.8314255,-0.1767249) (0.8343831,-0.1621876) (0.8370866,-0.1476010) (0.8395351,-0.1329693) (0.8417279,-0.1182971) (0.8436642,-0.1035889) (0.8453436,-0.0888492) (0.8467655,-0.0740824) (0.8479294,-0.0592930) (0.8488351,-0.0444856) (0.8494822,-0.0296646) (0.8498705,-0.0148345) (0.8500000,0.0000000) (0.8498705,0.0148345) (0.8494822,0.0296646) (0.8488351,0.0444856) (0.8479294,0.0592930) (0.8467655,0.0740824) (0.8453436,0.0888492) (0.8436642,0.1035889) (0.8417279,0.1182971) (0.8395351,0.1329693) (0.8370866,0.1476010) (0.8343831,0.1621876) (0.8314255,0.1767249) (0.8282146,0.1912084) (0.8247514,0.2056336) (0.8210370,0.2199962) (0.8170724,0.2342918) (0.8128590,0.2485159) (0.8083980,0.2626644) (0.8036908,0.2767329) (0.7987387,0.2907171) (0.7935434,0.3046128) (0.7881063,0.3184156) (0.7824291,0.3321215) (0.7765136,0.3457261) (0.7703616,0.3592255) (0.7639749,0.3726155) (0.7573555,0.3858919) (0.7505055,0.3990508) (0.7434268,0.4120882) (0.7361216,0.4250000) (0.7285922,0.4377824) (0.7208409,0.4504314) (0.7128700,0.4629432) (0.7046819,0.4753140) (0.6962792,0.4875400) (0.6876644,0.4996175) (0.6788402,0.5115428) (0.6698091,0.5233123) (0.6605741,0.5349223) (0.6511378,0.5463695) (0.6415031,0.5576502) (0.6316731,0.5687610) (0.6216506,0.5796986) (0.6114388,0.5904596) (0.6010408,0.6010408) (0.5904596,0.6114388) (0.5796986,0.6216506) (0.5687610,0.6316731) (0.5576502,0.6415031) (0.5463695,0.6511378)};
\draw[thick] plot coordinates {(0.8034845,-0.9575556) (0.8200738,-0.9433870) (0.8364133,-0.9289310) (0.8524980,-0.9141921) (0.8683230,-0.8991748) (0.8838835,-0.8838835) (0.8991748,-0.8683230) (0.9141921,-0.8524980) (0.9289310,-0.8364133) (0.9433870,-0.8200738) (0.9575556,-0.8034845) (0.9714325,-0.7866505) (0.9850134,-0.7695768) (0.9982944,-0.7522688) (1.0112712,-0.7347316) (1.0239401,-0.7169705) (1.0362970,-0.6989911) (1.0483382,-0.6807988) (1.0600601,-0.6623991) (1.0714591,-0.6437976) (1.0825318,-0.6250000) (1.0932746,-0.6060120) (1.1036845,-0.5868395) (1.1137582,-0.5674881) (1.1234926,-0.5479639) (1.1328847,-0.5282728) (1.1419318,-0.5084208) (1.1506311,-0.4884139) (1.1589798,-0.4682582) (1.1669755,-0.4479599) (1.1746158,-0.4275252) (1.1818982,-0.4069602) (1.1888206,-0.3862712) (1.1953809,-0.3654646) (1.2015771,-0.3445467) (1.2074073,-0.3235238) (1.2128697,-0.3024024) (1.2179626,-0.2811888) (1.2226845,-0.2598896) (1.2270340,-0.2385112) (1.2310097,-0.2170602) (1.2346104,-0.1955431) (1.2378351,-0.1739664) (1.2406827,-0.1523367) (1.2431524,-0.1306606) (1.2452434,-0.1089447) (1.2469551,-0.0871956) (1.2482869,-0.0654199) (1.2492385,-0.0436244) (1.2498096,-0.0218155) (1.2500000,0.0000000) (1.2498096,0.0218155) (1.2492385,0.0436244) (1.2482869,0.0654199) (1.2469551,0.0871956) (1.2452434,0.1089447) (1.2431524,0.1306606) (1.2406827,0.1523367) (1.2378351,0.1739664) (1.2346104,0.1955431) (1.2310097,0.2170602) (1.2270340,0.2385112) (1.2226845,0.2598896) (1.2179626,0.2811888) (1.2128697,0.3024024) (1.2074073,0.3235238) (1.2015771,0.3445467) (1.1953809,0.3654646) (1.1888206,0.3862712) (1.1818982,0.4069602) (1.1746158,0.4275252) (1.1669755,0.4479599) (1.1589798,0.4682582) (1.1506311,0.4884139) (1.1419318,0.5084208) (1.1328847,0.5282728) (1.1234926,0.5479639) (1.1137582,0.5674881) (1.1036845,0.5868395) (1.0932746,0.6060120) (1.0825318,0.6250000) (1.0714591,0.6437976) (1.0600601,0.6623991) (1.0483382,0.6807988) (1.0362970,0.6989911) (1.0239401,0.7169705) (1.0112712,0.7347316) (0.9982944,0.7522688) (0.9850134,0.7695768) (0.9714325,0.7866505) (0.9575556,0.8034845) (0.9433870,0.8200738) (0.9289310,0.8364133) (0.9141921,0.8524980) (0.8991748,0.8683230) (0.8838835,0.8838835) (0.8683230,0.8991748) (0.8524980,0.9141921) (0.8364133,0.9289310) (0.8200738,0.9433870) (0.8034845,0.9575556)};
\draw[DeepBlue,thick,->] (0.0000000,0.0000000)--(1.6500000,0.0000000);
\node[below,fill=none,inner sep=2pt] at (0.8000000,-1.1500000) {同心波面与法线};
\end{scope}
\end{tikzpicture}
\figcap{观测点 $O$ 在地表；按到眼光线的切线反向看，太阳视方向高于真实方向。左侧为同心波面。}\label{ch27:fig:s42-19-1}
\end{minipage}\end{center}
}

\newcommand{\ReflectionPanel}{%
\begin{center}\begin{minipage}{\linewidth}\centering
\begin{tikzpicture}[>=Latex,font=\small,every node/.style={fill=none}]
\begin{scope}[shift={(-4.1000000,0.0000000)}]
\draw[DeepBlue,thick,->] (-1.4000000,1.4000000)--(0.0000000,0.0000000);
\draw[DeepBlue,thick,->] (0.0000000,0.0000000)--(1.4000000,1.4000000);
\draw[thick] (-1.8000000,0.0000000)--(1.8000000,0.0000000);
\draw[gray,dashed,thin] (0.0000000,-0.2000000)--(0.0000000,1.8000000);
\node[above,fill=none,inner sep=2pt] at (-0.4000000,0.8000000) {$i$};
\node[above,fill=none,inner sep=2pt] at (0.4000000,0.8000000) {$r$};
\node[below,fill=none,inner sep=2pt] at (0.0000000,-0.3500000) {反射角等于入射角};
\end{scope}
\begin{scope}[shift={(0.3000000,0.0000000)}]
\draw[thick] (-1.9336619,-0.4762222)--(0.1336619,0.2762222);
\draw[thick] (-0.5547819,1.4332399)--(1.5125419,0.6807955);
\draw[DeepBlue,thick,->] (-2.4000000,-0.1000000)--(-0.9000000,-0.1000000);
\draw[DeepBlue,thick,->] (-0.9000000,-0.1000000)--(0.4788800,1.0570177);
\draw[DeepBlue,thick,->] (0.4788800,1.0570177)--(0.7219874,-0.3217132);
\node[below,fill=none,inner sep=2pt] at (0.0000000,-0.7000000) {两镜夹角 $40^\circ$};
\node[below,fill=none,inner sep=2pt] at (0.0000000,-1.1000000) {方向转过 $80^\circ$};
\end{scope}
\begin{scope}[shift={(4.7000000,0.0000000)}]
\draw[thick] (-1.2000000,1.3000000)--(-1.2000000,-1.2000000)--(1.6000000,-1.2000000);
\draw[DeepBlue,thick,->] (1.5000000,1.0000000)--(-1.2000000,-0.3500000);
\draw[DeepBlue,thick,->] (-1.2000000,-0.3500000)--(0.5000000,-1.2000000);
\draw[DeepBlue,thick,->] (0.5000000,-1.2000000)--(2.0000000,-0.4500000);
\node[below,fill=none,inner sep=2pt] at (0.1500000,-1.5000000) {直角反射器};
\end{scope}
\end{tikzpicture}
\figcap{单镜、双镜与直角反射器。各段在镜面交点处等角反射。}\label{ch27:fig:s42-19-3}
\end{minipage}\end{center}
}

\newcommand{\RefractionPrismPanel}{%
\begin{center}\begin{minipage}{\linewidth}\centering
\begin{tikzpicture}[>=Latex,font=\small,every node/.style={fill=none}]
\begin{scope}[shift={(-3.2000000,0.0000000)}]
\draw[thick] (-2.0000000,0.0000000)--(2.0000000,0.0000000);
\draw[gray,dashed,thin] (0.0000000,-1.8000000)--(0.0000000,1.8000000);
\draw[DeepBlue,thick,->] (-1.5000000,1.5000000)--(0.0000000,0.0000000);
\draw[DeepBlue,thick,->] (0.0000000,0.0000000)--(0.8017840,-1.5000000);
\node[above,fill=none,inner sep=2pt] at (-1.6000000,0.4000000) {$n_1=1$};
\node[above,fill=none,inner sep=2pt] at (-1.5000000,-1.1000000) {$n_2=1.5$};
\node[above,fill=none,inner sep=2pt] at (-0.3500000,0.9000000) {$i$};
\node[right,fill=none,inner sep=2pt] at (0.7500000,-0.8500000) {$r$};
\draw[gray,dashed,thin] (0.7200000,-0.8500000)--(0.2000000,-0.7000000);
\end{scope}
\begin{scope}[shift={(3.0000000,0.0000000)}]
\draw[thick] (-1.5000000,-1.0000000)--(0.0000000,1.5980760)--(1.5000000,-1.0000000)--(-1.5000000,-1.0000000);
\draw[DeepBlue,thick,->] (-2.4167961,-0.3128935)--(-0.8000000,0.2124354);
\draw[gray,dashed,thin] (-1.2763140,0.4874354)--(-0.1937822,-0.1375646);
\draw[red!65!black,thick,->] (-0.8000000,0.2124354)--(0.8048823,0.2039789);
\draw[red!65!black,thick,->] (0.8048823,0.2039789)--(2.1743453,-0.2725409);
\draw[gray,dashed,thin] (0.3718696,-0.0460211)--(1.3677988,0.5289789);
\draw[Blue,thick,->] (-0.8000000,0.2124354)--(0.8153120,0.1859141);
\draw[Blue,thick,->] (0.8153120,0.1859141)--(2.1609251,-0.3542946);
\node[above,fill=none,inner sep=2pt] at (0.0000000,1.0000000) {$A$};
\node[below,fill=none,inner sep=2pt] at (0.0000000,-1.3000000) {各颜色只在界面改变方向};
\end{scope}
\end{tikzpicture}
\figcap{平界面折射与棱镜色散；虚线为界面法线。}\label{ch27:fig:s42-19-6}
\end{minipage}\end{center}
}

\newcommand{\PrismDeviationDiagram}{%
\begin{center}\begin{minipage}{\linewidth}\centering
\begin{tikzpicture}[>=Latex,font=\small,every node/.style={fill=none}]
\draw[gray,thin,->] (0.0000000,0.0000000)--(6.1000000,0.0000000);
\draw[gray,thin,->] (0.0000000,0.0000000)--(0.0000000,3.3000000);
\draw[DeepBlue,thick] plot coordinates {(0.6601178,2.2650478) (0.6799715,1.8798761) (0.6998524,1.7231931) (0.7197610,1.6081283) (0.7396977,1.5145945) (0.7596629,1.4348059) (0.7796570,1.3647702) (0.7996804,1.3021169) (0.8197336,1.2453026) (0.8398170,1.1932560) (0.8599310,1.1451978) (0.8800762,1.1005403) (0.9002529,1.0588274) (0.9204616,1.0196967) (0.9407028,0.9828547) (0.9609769,0.9480600) (0.9812845,0.9151108) (1.0016261,0.8838370) (1.0220022,0.8540931) (1.0424132,0.8257539) (1.0628597,0.7987107) (1.0833423,0.7728683) (1.1038615,0.7481429) (1.1244178,0.7244604) (1.1450117,0.7017545) (1.1656440,0.6799662) (1.1863151,0.6590424) (1.2070256,0.6389350) (1.2277763,0.6196005) (1.2485675,0.6009996) (1.2694001,0.5830962) (1.2902746,0.5658573) (1.3111918,0.5492529) (1.3321521,0.5332550) (1.3531565,0.5178383) (1.3742054,0.5029790) (1.3952997,0.4886554) (1.4164401,0.4748471) (1.4376272,0.4615356) (1.4588620,0.4487032) (1.4801450,0.4363337) (1.5014772,0.4244120) (1.5228592,0.4129239) (1.5442920,0.4018562) (1.5657763,0.3911966) (1.5873131,0.3809334) (1.6089031,0.3710558) (1.6305473,0.3615536) (1.6522466,0.3524172) (1.6740020,0.3436378) (1.6958143,0.3352069) (1.7176845,0.3271166) (1.7396137,0.3193594) (1.7616029,0.3119284) (1.7836531,0.3048170) (1.8057654,0.2980190) (1.8279409,0.2915288) (1.8501807,0.2853407) (1.8724860,0.2794497) (1.8948580,0.2738511) (1.9172978,0.2685403) (1.9398067,0.2635132) (1.9623860,0.2587658) (1.9850371,0.2542945) (2.0077611,0.2500959) (2.0305596,0.2461668) (2.0534339,0.2425043) (2.0763854,0.2391057) (2.0994158,0.2359685) (2.1225265,0.2330903) (2.1457191,0.2304693) (2.1689952,0.2281034) (2.1923565,0.2259909) (2.2158047,0.2241305) (2.2393416,0.2225207) (2.2629690,0.2211605) (2.2866888,0.2200489) (2.3105029,0.2191851) (2.3344134,0.2185686) (2.3584222,0.2181988) (2.3825315,0.2180756) (2.4067435,0.2181988) (2.4310604,0.2185686) (2.4554847,0.2191851) (2.4800186,0.2200489) (2.5046648,0.2211605) (2.5294257,0.2225207) (2.5543040,0.2241305) (2.5793026,0.2259909) (2.6044243,0.2281034) (2.6296720,0.2304693) (2.6550488,0.2330903) (2.6805579,0.2359685) (2.7062026,0.2391057) (2.7319864,0.2425043) (2.7579128,0.2461668) (2.7839855,0.2500959) (2.8102084,0.2542945) (2.8365855,0.2587658) (2.8631209,0.2635132) (2.8898191,0.2685403) (2.9166846,0.2738511) (2.9437221,0.2794497) (2.9709365,0.2853407) (2.9983330,0.2915288) (3.0259171,0.2980190) (3.0536944,0.3048170) (3.0816707,0.3119284) (3.1098524,0.3193594) (3.1382460,0.3271166) (3.1668582,0.3352069) (3.1956962,0.3436378) (3.2247677,0.3524172) (3.2540806,0.3615536) (3.2836434,0.3710558) (3.3134647,0.3809334) (3.3435542,0.3911966) (3.3739215,0.4018562) (3.4045774,0.4129239) (3.4355329,0.4244120) (3.4667998,0.4363337) (3.4983907,0.4487032) (3.5303191,0.4615356) (3.5625992,0.4748471) (3.5952464,0.4886554) (3.6282771,0.5029790) (3.6617088,0.5178383) (3.6955604,0.5332550) (3.7298523,0.5492529) (3.7646065,0.5658573) (3.7998467,0.5830962) (3.8355988,0.6009996) (3.8718909,0.6196005) (3.9087535,0.6389350) (3.9462202,0.6590424) (3.9843279,0.6799662) (4.0231170,0.7017545) (4.0626325,0.7244604) (4.1029243,0.7481429) (4.1440479,0.7728683) (4.1860658,0.7987107) (4.2290484,0.8257539) (4.2730754,0.8540931) (4.3182380,0.8838370) (4.3646411,0.9151108) (4.4124064,0.9480600) (4.4616762,0.9828547) (4.5126190,1.0196967) (4.5654366,1.0588274) (4.6203741,1.1005403) (4.6777338,1.1451978) (4.7378963,1.1932560) (4.8013519,1.2453026) (4.8687504,1.3021169) (4.9409849,1.3647702) (5.0193421,1.4348059) (5.1057977,1.5145945) (5.2036792,1.6081283) (5.3194752,1.7231931) (5.4699253,1.8798761) (5.8107553,2.2650478)};
\fill (2.3825315,0.2180756) circle (1.2pt);
\node[below,fill=none,inner sep=2pt] at (2.3825315,0.2180756) {};
\draw[gray,dashed,thin] (2.3825315,0.2180756)--(2.4000000,1.1500000);
\node[above,fill=none,inner sep=2pt] at (2.4000000,1.2000000) {对称光路};
\node[below,fill=none,inner sep=2pt] at (0.8333333,-0.0800000) {$30^\circ$};
\node[below,fill=none,inner sep=2pt] at (2.5000000,-0.0800000) {$50^\circ$};
\node[below,fill=none,inner sep=2pt] at (4.1666667,-0.0800000) {$70^\circ$};
\node[below,fill=none,inner sep=2pt] at (5.8333333,-0.0800000) {$90^\circ$};
\node[left,fill=none,inner sep=2pt] at (-0.0500000,0.5000000) {$40^\circ$};
\node[left,fill=none,inner sep=2pt] at (-0.0500000,1.5000000) {$50^\circ$};
\node[left,fill=none,inner sep=2pt] at (-0.0500000,2.5000000) {$60^\circ$};
\node[right,fill=none,inner sep=2pt] at (6.1000000,0.0000000) {$i_1$};
\node[above,fill=none,inner sep=2pt] at (0.0000000,3.3000000) {$\delta$};
\node[above,fill=none,inner sep=2pt] at (4.0000000,2.8000000) {$A=60^\circ,\ n=1.5$};
\end{tikzpicture}
\figcap{偏向角随入射角的变化；最低点是两内部折射角相等的光路。\label{ch27:fig:deviation}}
\end{minipage}\end{center}
}

\newcommand{\FiberFermatPanel}{%
\begin{center}\begin{minipage}{\linewidth}\centering
\begin{tikzpicture}[>=Latex,font=\small,every node/.style={fill=none}]
\begin{scope}[shift={(-3.7000000,0.0000000)}]
\draw[thick] (-2.6000000,-0.6500000)--(2.6000000,-0.6500000);
\draw[thick] (-2.6000000,0.6500000)--(2.6000000,0.6500000);
\draw[gray,dashed,thin] (-2.6000000,-0.9000000)--(2.6000000,-0.9000000);
\draw[gray,dashed,thin] (-2.6000000,0.9000000)--(2.6000000,0.9000000);
\draw[DeepBlue,thick,->] (-2.6000000,0.0000000)--(-1.9500000,0.6500000);
\draw[DeepBlue,thick,->] (-1.9500000,0.6500000)--(-0.6500000,-0.6500000);
\draw[DeepBlue,thick,->] (-0.6500000,-0.6500000)--(0.6500000,0.6500000);
\draw[DeepBlue,thick,->] (0.6500000,0.6500000)--(1.9500000,-0.6500000);
\draw[DeepBlue,thick,->] (1.9500000,-0.6500000)--(2.6000000,0.0000000);
\node[above,fill=none,inner sep=2pt] at (0.0000000,1.0500000) {包层 $n_2<n_1$};
\node[below,fill=none,inner sep=2pt] at (0.0000000,-1.0500000) {芯层内，反射点均在边界};
\end{scope}
\begin{scope}[shift={(3.0000000,0.0000000)}]
\draw[thick] plot coordinates {(0.9844742,-1.0500000) (0.9383572,-1.0290000) (0.8944127,-1.0080000) (0.8524231,-0.9870000) (0.8122095,-0.9660000) (0.7736233,-0.9450000) (0.7365391,-0.9240000) (0.7008507,-0.9030000) (0.6664670,-0.8820000) (0.6333096,-0.8610000) (0.6013107,-0.8400000) (0.5704111,-0.8190000) (0.5405591,-0.7980000) (0.5117096,-0.7770000) (0.4838228,-0.7560000) (0.4568639,-0.7350000) (0.4308022,-0.7140000) (0.4056108,-0.6930000) (0.3812660,-0.6720000) (0.3577471,-0.6510000) (0.3350360,-0.6300000) (0.3131170,-0.6090000) (0.2919765,-0.5880000) (0.2716029,-0.5670000) (0.2519862,-0.5460000) (0.2331184,-0.5250000) (0.2149924,-0.5040000) (0.1976029,-0.4830000) (0.1809455,-0.4620000) (0.1650171,-0.4410000) (0.1498155,-0.4200000) (0.1353393,-0.3990000) (0.1215882,-0.3780000) (0.1085622,-0.3570000) (0.0962623,-0.3360000) (0.0846897,-0.3150000) (0.0738464,-0.2940000) (0.0637344,-0.2730000) (0.0543561,-0.2520000) (0.0457143,-0.2310000) (0.0378117,-0.2100000) (0.0306511,-0.1890000) (0.0242354,-0.1680000) (0.0185671,-0.1470000) (0.0136488,-0.1260000) (0.0094830,-0.1050000) (0.0060716,-0.0840000) (0.0034163,-0.0630000) (0.0015187,-0.0420000) (0.0003797,-0.0210000) (0.0000000,0.0000000) (0.0003797,0.0210000) (0.0015187,0.0420000) (0.0034163,0.0630000) (0.0060716,0.0840000) (0.0094830,0.1050000) (0.0136488,0.1260000) (0.0185671,0.1470000) (0.0242354,0.1680000) (0.0306511,0.1890000) (0.0378117,0.2100000) (0.0457143,0.2310000) (0.0543561,0.2520000) (0.0637344,0.2730000) (0.0738464,0.2940000) (0.0846897,0.3150000) (0.0962623,0.3360000) (0.1085622,0.3570000) (0.1215882,0.3780000) (0.1353393,0.3990000) (0.1498155,0.4200000) (0.1650171,0.4410000) (0.1809455,0.4620000) (0.1976029,0.4830000) (0.2149924,0.5040000) (0.2331184,0.5250000) (0.2519862,0.5460000) (0.2716029,0.5670000) (0.2919765,0.5880000) (0.3131170,0.6090000) (0.3350360,0.6300000) (0.3577471,0.6510000) (0.3812660,0.6720000) (0.4056108,0.6930000) (0.4308022,0.7140000) (0.4568639,0.7350000) (0.4838228,0.7560000) (0.5117096,0.7770000) (0.5405591,0.7980000) (0.5704111,0.8190000) (0.6013107,0.8400000) (0.6333096,0.8610000) (0.6664670,0.8820000) (0.7008507,0.9030000) (0.7365391,0.9240000) (0.7736233,0.9450000) (0.8122095,0.9660000) (0.8524231,0.9870000) (0.8944127,1.0080000) (0.9383572,1.0290000) (0.9844742,1.0500000)};
\draw[gray,thin,->] (-2.3000000,0.0000000)--(3.4000000,0.0000000);
\fill (-2.0000000,0.0000000) circle (1.2pt);
\node[below,fill=none,inner sep=2pt] at (-2.0000000,0.0000000) {$O$};
\fill (3.0000000,0.0000000) circle (1.2pt);
\node[below,fill=none,inner sep=2pt] at (3.0000000,0.0000000) {$I$};
\draw[DeepBlue,thick,->] (-2.0000000,0.0000000)--(0.6164078,0.8500000);
\draw[DeepBlue,thick,->] (0.6164078,0.8500000)--(3.0000000,0.0000000);
\fill (0.6164078,0.8500000) circle (1.2pt);
\node[above,fill=none,inner sep=2pt] at (0.6164078,0.8500000) {$P$};
\node[below,fill=none,inner sep=2pt] at (0.0000000,-1.3500000) {$OP+nPI=s_o+ns_i$};
\end{scope}
\end{tikzpicture}
\figcap{左为平直芯包界面导光；右为等光程非球面，离轴点 $P$ 也连接同一物、像点。}\label{ch27:fig:s42-19-10}
\end{minipage}\end{center}
}

\newcommand{\SphericalMirrorRotatingDiagram}{%
\begin{center}\begin{minipage}{\linewidth}\centering
\begin{tikzpicture}[>=Latex,font=\small,every node/.style={fill=none}]
\begin{scope}[shift={(-3.0000000,0.0000000)}]
\draw[thick] plot coordinates {(2.3640323,-1.8469844) (2.3883229,-1.8154652) (2.4121933,-1.7836265) (2.4356393,-1.7514739) (2.4586568,-1.7190133) (2.4812417,-1.6862501) (2.5033901,-1.6531903) (2.5250979,-1.6198396) (2.5463615,-1.5862039) (2.5671771,-1.5522892) (2.5875410,-1.5181013) (2.6074496,-1.4836463) (2.6268995,-1.4489303) (2.6458872,-1.4139593) (2.6644093,-1.3787396) (2.6824627,-1.3432773) (2.7000441,-1.3075786) (2.7171505,-1.2716499) (2.7337787,-1.2354974) (2.7499260,-1.1991276) (2.7655895,-1.1625468) (2.7807663,-1.1257614) (2.7954539,-1.0887780) (2.8096497,-1.0516030) (2.8233511,-1.0142429) (2.8365557,-0.9767045) (2.8492613,-0.9389941) (2.8614656,-0.9011186) (2.8731664,-0.8630845) (2.8843617,-0.8248986) (2.8950495,-0.7865675) (2.9052279,-0.7480981) (2.9148952,-0.7094970) (2.9240496,-0.6707711) (2.9326896,-0.6319271) (2.9408135,-0.5929720) (2.9484201,-0.5539126) (2.9555078,-0.5147557) (2.9620756,-0.4755082) (2.9681222,-0.4361771) (2.9736466,-0.3967692) (2.9786478,-0.3572915) (2.9831249,-0.3177509) (2.9870772,-0.2781545) (2.9905039,-0.2385091) (2.9934044,-0.1988217) (2.9957783,-0.1590994) (2.9976250,-0.1193490) (2.9989444,-0.0795777) (2.9997361,-0.0397923) (3.0000000,0.0000000) (2.9997361,0.0397923) (2.9989444,0.0795777) (2.9976250,0.1193490) (2.9957783,0.1590994) (2.9934044,0.1988217) (2.9905039,0.2385091) (2.9870772,0.2781545) (2.9831249,0.3177509) (2.9786478,0.3572915) (2.9736466,0.3967692) (2.9681222,0.4361771) (2.9620756,0.4755082) (2.9555078,0.5147557) (2.9484201,0.5539126) (2.9408135,0.5929720) (2.9326896,0.6319271) (2.9240496,0.6707711) (2.9148952,0.7094970) (2.9052279,0.7480981) (2.8950495,0.7865675) (2.8843617,0.8248986) (2.8731664,0.8630845) (2.8614656,0.9011186) (2.8492613,0.9389941) (2.8365557,0.9767045) (2.8233511,1.0142429) (2.8096497,1.0516030) (2.7954539,1.0887780) (2.7807663,1.1257614) (2.7655895,1.1625468) (2.7499260,1.1991276) (2.7337787,1.2354974) (2.7171505,1.2716499) (2.7000441,1.3075786) (2.6824627,1.3432773) (2.6644093,1.3787396) (2.6458872,1.4139593) (2.6268995,1.4489303) (2.6074496,1.4836463) (2.5875410,1.5181013) (2.5671771,1.5522892) (2.5463615,1.5862039) (2.5250979,1.6198396) (2.5033901,1.6531903) (2.4812417,1.6862501) (2.4586568,1.7190133) (2.4356393,1.7514739) (2.4121933,1.7836265) (2.3883229,1.8154652) (2.3640323,1.8469844)};
\draw[gray,thin,->] (-0.7000000,0.0000000)--(3.6000000,0.0000000);
\fill (0.0000000,0.0000000) circle (1.2pt);
\node[below,fill=none,inner sep=2pt] at (0.0000000,0.0000000) {$C$};
\fill (3.0000000,0.0000000) circle (1.2pt);
\node[below right,fill=none,inner sep=2pt] at (3.0000000,0.0000000) {$V$};
\fill (1.6550669,0.0000000) circle (1.2pt);
\node[below,fill=none,inner sep=2pt] at (1.6550669,0.0000000) {$M$};
\draw[DeepBlue,thick,->] (-1.0000000,1.2678548)--(2.7189234,1.2678548);
\draw[DeepBlue,thick,->] (2.7189234,1.2678548)--(1.6550669,0.0000000);
\draw[gray,dashed,thin] (0.0000000,0.0000000)--(2.7189234,1.2678548);
\fill (2.7189234,1.2678548) circle (1.2pt);
\node[below,fill=none,inner sep=2pt] at (2.7189234,1.2678548) {};
\draw[gray,dashed,thin] (2.7189234,1.2678548)--(2.1689234,1.7178548);
\node[above left,fill=none,inner sep=2pt] at (2.1189234,1.7678548) {$P$};
\end{scope}
\begin{scope}[shift={(4.0000000,0.0000000)}]
\draw[thick] (0.0000000,-0.3000000)--(0.0000000,2.6000000);
\draw[thick] (0.0000000,0.0000000)--(-0.9500000,2.3513325);
\fill (0.0000000,0.0000000) circle (1.2pt);
\node[below,fill=none,inner sep=2pt] at (0.0000000,0.0000000) {$O$};
\fill (0.0000000,1.0000000) circle (1.2pt);
\node[right,fill=none,inner sep=2pt] at (0.0000000,1.0000000) {$A$};
\fill (-0.4040262,1.0000000) circle (1.2pt);
\node[left,fill=none,inner sep=2pt] at (-0.4040262,1.0000000) {$C$};
\fill (0.0000000,1.3901636) circle (1.2pt);
\node[right,fill=none,inner sep=2pt] at (0.0000000,1.3901636) {$D$};
\draw[DeepBlue,thick,->] (0.0000000,1.0000000)--(-0.4040262,1.0000000);
\draw[DeepBlue,thick,->] (-0.4040262,1.0000000)--(0.0000000,1.3901636);
\draw[thick] plot coordinates {(0.0000000,0.4200000) (-0.0016127,0.4199969) (-0.0032253,0.4199876) (-0.0048379,0.4199721) (-0.0064505,0.4199505) (-0.0080629,0.4199226) (-0.0096752,0.4198885) (-0.0112874,0.4198483) (-0.0128994,0.4198019) (-0.0145113,0.4197492) (-0.0161229,0.4196904) (-0.0177343,0.4196254) (-0.0193454,0.4195542) (-0.0209562,0.4194769) (-0.0225667,0.4193933) (-0.0241769,0.4193036) (-0.0257867,0.4192076) (-0.0273962,0.4191055) (-0.0290052,0.4189973) (-0.0306138,0.4188828) (-0.0322220,0.4187622) (-0.0338297,0.4186353) (-0.0354369,0.4185024) (-0.0370435,0.4183632) (-0.0386497,0.4182179) (-0.0402552,0.4180664) (-0.0418602,0.4179088) (-0.0434645,0.4177449) (-0.0450682,0.4175750) (-0.0466713,0.4173988) (-0.0482736,0.4172166) (-0.0498752,0.4170281) (-0.0514761,0.4168335) (-0.0530763,0.4166328) (-0.0546756,0.4164260) (-0.0562742,0.4162129) (-0.0578719,0.4159938) (-0.0594688,0.4157685) (-0.0610648,0.4155371) (-0.0626599,0.4152996) (-0.0642541,0.4150559) (-0.0658473,0.4148061) (-0.0674395,0.4145503) (-0.0690308,0.4142882) (-0.0706210,0.4140201) (-0.0722102,0.4137459) (-0.0737984,0.4134656) (-0.0753854,0.4131792) (-0.0769713,0.4128867) (-0.0785561,0.4125881) (-0.0801398,0.4122834) (-0.0817222,0.4119727) (-0.0833035,0.4116558) (-0.0848835,0.4113329) (-0.0864623,0.4110040) (-0.0880398,0.4106690) (-0.0896160,0.4103279) (-0.0911909,0.4099808) (-0.0927644,0.4096276) (-0.0943366,0.4092684) (-0.0959074,0.4089031) (-0.0974767,0.4085319) (-0.0990447,0.4081546) (-0.1006111,0.4077713) (-0.1021761,0.4073819) (-0.1037396,0.4069866) (-0.1053015,0.4065853) (-0.1068619,0.4061780) (-0.1084207,0.4057646) (-0.1099780,0.4053453) (-0.1115336,0.4049201) (-0.1130875,0.4044888) (-0.1146398,0.4040516) (-0.1161904,0.4036085) (-0.1177393,0.4031593) (-0.1192864,0.4027043) (-0.1208318,0.4022433) (-0.1223754,0.4017764) (-0.1239172,0.4013035) (-0.1254572,0.4008248) (-0.1269954,0.4003401) (-0.1285316,0.3998495) (-0.1300660,0.3993530) (-0.1315984,0.3988507) (-0.1331289,0.3983424) (-0.1346575,0.3978283) (-0.1361840,0.3973083) (-0.1377086,0.3967825) (-0.1392311,0.3962508) (-0.1407515,0.3957133) (-0.1422699,0.3951699) (-0.1437862,0.3946207) (-0.1453004,0.3940657) (-0.1468124,0.3935049) (-0.1483223,0.3929383) (-0.1498300,0.3923659) (-0.1513354,0.3917877) (-0.1528387,0.3912037) (-0.1543396,0.3906140) (-0.1558384,0.3900185) (-0.1573348,0.3894172)};
\node[left,fill=none,inner sep=2pt] at (-0.8000000,0.3500000) {$\theta$};
\draw[gray,dashed,thin] (-0.7500000,0.4000000)--(-0.0800000,0.4200000);
\draw[gray,<->] (0.6000000,0.0000000)--(0.6000000,1.0000000);
\node[right,fill=none,inner sep=2pt] at (0.6000000,0.5000000) {$d$};
\end{scope}
\end{tikzpicture}
\figcap{左侧球面镜的 $P$ 在圆周上，$CP$ 为法线；右侧转镜的转轴 $O$ 位于墙上。}\label{ch27:fig:s42-19-11}
\end{minipage}\end{center}
}

\newcommand{\CylinderDiagram}{%
\begin{center}\begin{minipage}{\linewidth}\centering
\begin{tikzpicture}[>=Latex,font=\small,every node/.style={fill=none}]
\draw[thick] (0.0000000,0.0000000) circle (2);
\draw[thick] plot coordinates {(1.8131983,-0.9238707) (1.8218249,-0.9067409) (1.8302897,-0.8895306) (1.8385919,-0.8722413) (1.8467308,-0.8548745) (1.8547057,-0.8374317) (1.8625158,-0.8199146) (1.8701605,-0.8023246) (1.8776390,-0.7846634) (1.8849508,-0.7669325) (1.8920951,-0.7491335) (1.8990714,-0.7312679) (1.9058790,-0.7133373) (1.9125173,-0.6953434) (1.9189857,-0.6772878) (1.9252837,-0.6591719) (1.9314107,-0.6409976) (1.9373660,-0.6227663) (1.9431493,-0.6044796) (1.9487600,-0.5861393) (1.9541977,-0.5677469) (1.9594617,-0.5493041) (1.9645516,-0.5308125) (1.9694671,-0.5122737) (1.9742076,-0.4936894) (1.9787728,-0.4750613) (1.9831622,-0.4563910) (1.9873754,-0.4376801) (1.9914121,-0.4189304) (1.9952720,-0.4001435) (1.9989546,-0.3813210) (2.0024596,-0.3624646) (2.0057868,-0.3435761) (2.0089357,-0.3246570) (2.0119063,-0.3057091) (2.0146981,-0.2867340) (2.0173110,-0.2677335) (2.0197447,-0.2487092) (2.0219990,-0.2296627) (2.0240737,-0.2105959) (2.0259686,-0.1915104) (2.0276835,-0.1724079) (2.0292184,-0.1532900) (2.0305729,-0.1341586) (2.0317472,-0.1150152) (2.0327409,-0.0958616) (2.0335541,-0.0766995) (2.0341866,-0.0575306) (2.0346385,-0.0383566) (2.0349096,-0.0191791) (2.0350000,0.0000000) (2.0349096,0.0191791) (2.0346385,0.0383566) (2.0341866,0.0575306) (2.0335541,0.0766995) (2.0327409,0.0958616) (2.0317472,0.1150152) (2.0305729,0.1341586) (2.0292184,0.1532900) (2.0276835,0.1724079) (2.0259686,0.1915104) (2.0240737,0.2105959) (2.0219990,0.2296627) (2.0197447,0.2487092) (2.0173110,0.2677335) (2.0146981,0.2867340) (2.0119063,0.3057091) (2.0089357,0.3246570) (2.0057868,0.3435761) (2.0024596,0.3624646) (1.9989546,0.3813210) (1.9952720,0.4001435) (1.9914121,0.4189304) (1.9873754,0.4376801) (1.9831622,0.4563910) (1.9787728,0.4750613) (1.9742076,0.4936894) (1.9694671,0.5122737) (1.9645516,0.5308125) (1.9594617,0.5493041) (1.9541977,0.5677469) (1.9487600,0.5861393) (1.9431493,0.6044796) (1.9373660,0.6227663) (1.9314107,0.6409976) (1.9252837,0.6591719) (1.9189857,0.6772878) (1.9125173,0.6953434) (1.9058790,0.7133373) (1.8990714,0.7312679) (1.8920951,0.7491335) (1.8849508,0.7669325) (1.8776390,0.7846634) (1.8701605,0.8023246) (1.8625158,0.8199146) (1.8547057,0.8374317) (1.8467308,0.8548745) (1.8385919,0.8722413) (1.8302897,0.8895306) (1.8218249,0.9067409) (1.8131983,0.9238707)};
\draw[gray,thin,->] (-2.4000000,0.0000000)--(2.5000000,0.0000000);
\draw[DeepBlue,thick,->] (-3.2000000,1.0000000)--(-1.7320508,1.0000000);
\draw[DeepBlue,thick,->] (-1.7320508,1.0000000)--(2.0000000,0.0000000);
\draw[DeepBlue,thick,->] (2.0000000,0.0000000)--(-1.7320508,-1.0000000);
\draw[DeepBlue,thick,->] (-1.7320508,-1.0000000)--(-3.2000000,-1.0000000);
\draw[gray,dashed,thin] (-1.7320508,1.0000000)--(0.0000000,0.0000000);
\draw[gray,dashed,thin] (-1.7320508,-1.0000000)--(0.0000000,0.0000000);
\fill (-1.7320508,1.0000000) circle (1.2pt);
\node[above left,fill=none,inner sep=2pt] at (-1.7320508,1.0000000) {$A$};
\fill (-1.7320508,-1.0000000) circle (1.2pt);
\node[below left,fill=none,inner sep=2pt] at (-1.7320508,-1.0000000) {$B$};
\fill (2.0000000,0.0000000) circle (1.2pt);
\node[right,fill=none,inner sep=2pt] at (2.0000000,0.0000000) {$P$};
\fill (0.0000000,0.0000000) circle (1.2pt);
\node[below,fill=none,inner sep=2pt] at (0.0000000,0.0000000) {$C$};
\draw[gray,<->] (-3.5000000,-1.0000000)--(-3.5000000,1.0000000);
\node[left,fill=none,inner sep=2pt] at (-3.5000000,0.0000000) {$d$};
\node[below,fill=none,inner sep=2pt] at (1.5000000,-2.2000000) {右侧镀银；两外部光线反向平行};
\end{tikzpicture}
\figcap{$d$ 是入射与出射两平行线的间距；圆内两段只在反射面 $P$ 处折转。}\label{ch27:fig:s42-19-12}
\end{minipage}\end{center}
}

\newcommand{\AbsorbingShellDiagram}{%
% Independent panel: original physical paths and local coordinates retained.
\begin{center}\begin{minipage}{\linewidth}\centering
\begin{tikzpicture}[>=Latex,font=\small,every node/.style={fill=none}]
\draw[thick] (0.0000000,0.0000000) circle (2);
\draw[very thick] (0.0000000,0.0000000) circle (0.8);
\draw[DeepBlue,thick,->] (-3.0000000,1.2000000)--(-1.6000000,1.2000000);
\draw[DeepBlue,thick,->] (-1.6000000,1.2000000)--(0.1839273,0.7785697);
\draw[gray,dashed,thin] (-1.6000000,1.2000000)--(0.0000000,0.0000000);
\draw[gray,dashed,thin] (0.0000000,0.0000000)--(0.1839273,0.7785697);
\fill (0.1839273,0.7785697) circle (1.2pt);
\node[below,fill=none,inner sep=2pt] at (0.1839273,0.7785697) {};
\node[above right,fill=none,inner sep=2pt] at (0.4339273,1.0285697) {$T$};
\draw[gray,<->] (0.0000000,0.0000000)--(2.0000000,0.0000000);
\node[below,fill=none,inner sep=2pt] at (1.0000000,0.0000000) {$b$};
\node[above,fill=none,inner sep=2pt] at (-0.4000000,-0.2000000) {$a$};
\node[below,fill=none,inner sep=2pt] at (0.0000000,-2.2500000) {内球面为吸收层};
\end{tikzpicture}
\figcap{临界折射线在 $T$ 与半径为 $a$ 的吸收层相切，外球半径为 $b$。}\label{ch27:fig:s42-19-13}
\end{minipage}\end{center}
}

\newcommand{\GradientCircleDiagram}{%
% Independent panel: original physical paths and local coordinates retained.
\begin{center}\begin{minipage}{\linewidth}\centering
\begin{tikzpicture}[>=Latex,font=\small,every node/.style={fill=none}]
\draw[gray,thin,->] (-0.6000000,0.0000000)--(3.2000000,0.0000000);
\draw[gray,thin,->] (0.0000000,-0.3000000)--(0.0000000,3.2000000);
\node[right,fill=none,inner sep=2pt] at (3.2000000,0.0000000) {$x$};
\node[above,fill=none,inner sep=2pt] at (0.0000000,3.2000000) {$y$};
\draw[DeepBlue,thick,->] (-0.6000000,0.0000000)--(0.0000000,0.0000000);
\draw[DeepBlue,thick,->] plot coordinates {(0.0000000,0.0000000) (0.0347776,0.0002016) (0.0695505,0.0008063) (0.1043141,0.0018141) (0.1390637,0.0032249) (0.1737946,0.0050383) (0.2085021,0.0072543) (0.2431816,0.0098725) (0.2778284,0.0128925) (0.3124379,0.0163139) (0.3470054,0.0201364) (0.3815262,0.0243593) (0.4159958,0.0289821) (0.4504095,0.0340042) (0.4847626,0.0394249) (0.5190506,0.0452434) (0.5532689,0.0514591) (0.5874128,0.0580710) (0.6214777,0.0650783) (0.6554591,0.0724800) (0.6893525,0.0802752) (0.7231532,0.0884627) (0.7568567,0.0970415) (0.7904584,0.1060105) (0.8239540,0.1153683) (0.8573388,0.1251139) (0.8906084,0.1352458) (0.9237584,0.1457627) (0.9567841,0.1566632) (0.9896813,0.1679458) (1.0224455,0.1796090) (1.0550723,0.1916513) (1.0875573,0.2040710) (1.1198961,0.2168664) (1.1520844,0.2300358) (1.1841179,0.2435775) (1.2159923,0.2574897) (1.2477032,0.2717704) (1.2792465,0.2864178) (1.3106178,0.3014298) (1.3418130,0.3168045) (1.3728279,0.3325399) (1.4036583,0.3486337) (1.4343000,0.3650838) (1.4647490,0.3818880) (1.4950011,0.3990441) (1.5250523,0.4165497) (1.5548986,0.4344025) (1.5845359,0.4526001) (1.6139602,0.4711401) (1.6431677,0.4900199) (1.6721543,0.5092371) (1.7009162,0.5287889) (1.7294495,0.5486729) (1.7577504,0.5688863) (1.7858150,0.5894265) (1.8136397,0.6102906) (1.8412206,0.6314759) (1.8685541,0.6529794) (1.8956365,0.6747984) (1.9224641,0.6969299) (1.9490333,0.7193709) (1.9753406,0.7421184) (2.0013824,0.7651693) (2.0271553,0.7885206) (2.0526557,0.8121690) (2.0778803,0.8361115) (2.1028257,0.8603448) (2.1274884,0.8848657) (2.1518652,0.9096708) (2.1759528,0.9347568) (2.1997480,0.9601204) (2.2232476,0.9857581) (2.2464483,1.0116666) (2.2693472,1.0378422) (2.2919411,1.0642815) (2.3142270,1.0909810) (2.3362018,1.1179371) (2.3578627,1.1451460) (2.3792068,1.1726043) (2.4002310,1.2003081) (2.4209327,1.2282538) (2.4413091,1.2564376) (2.4613574,1.2848557) (2.4810748,1.3135043) (2.5004589,1.3423796) (2.5195069,1.3714776) (2.5382163,1.4007946) (2.5565846,1.4303264) (2.5746093,1.4600691) (2.5922880,1.4900189) (2.6096183,1.5201715) (2.6265979,1.5505231) (2.6432245,1.5810694) (2.6594959,1.6118064) (2.6754099,1.6427300) (2.6909643,1.6738360) (2.7061571,1.7051202) (2.7209862,1.7365784) (2.7354496,1.7682064) (2.7495454,1.8000000)};
\draw[gray,dashed,thin] (0.0000000,3.0000000)--(2.7495454,1.8000000);
\draw[gray,dashed,thin] (0.0000000,3.0000000)--(0.0000000,0.0000000);
\fill (0.0000000,3.0000000) circle (1.2pt);
\node[left,fill=none,inner sep=2pt] at (0.0000000,3.0000000) {$C$};
\node[below,fill=none,inner sep=2pt] at (0.8000000,2.1000000) {$R$};
\node[below,fill=none,inner sep=2pt] at (1.9000000,-0.3000000) {$x=R\sin\theta$};
\node[below,fill=none,inner sep=2pt] at (1.9000000,-0.7500000) {$y=R(1-\cos\theta)$};
\end{tikzpicture}
\figcap{光线从原点水平入射，沿圆心为 $C$ 的圆弧向上弯曲；圆弧切线确定传播方向。}\label{ch27:fig:s42-19-14}
\end{minipage}\end{center}
}

\newcommand{\FisherBayDiagram}{%
% Independent panel: original physical paths and local coordinates retained.
\begin{center}\begin{minipage}{\linewidth}\centering
\begin{tikzpicture}[>=Latex,font=\small,every node/.style={fill=none}]
\draw[thick] (-0.9576305,-3.5739256)--(0.0000000,0.0000000)--(0.9576305,-3.5739256);
\draw[gray,dashed,thin] (0.0000000,0.0000000)--(0.0000000,-3.6000000);
\draw[gray,dashed,thin] (0.0000000,0.0000000)--(-1.4974340,-2.4975371);
\draw[gray,dashed,thin] (-1.4974340,-2.4975371)--(-0.7246933,-2.7045923);
\fill (0.0000000,0.0000000) circle (1.2pt);
\node[above,fill=none,inner sep=2pt] at (0.0000000,0.0000000) {$D$};
\fill (-1.4974340,-2.4975371) circle (1.2pt);
\node[left,fill=none,inner sep=2pt] at (-1.4974340,-2.4975371) {$A$};
\fill (1.4974340,-2.4975371) circle (1.2pt);
\node[right,fill=none,inner sep=2pt] at (1.4974340,-2.4975371) {$B$};
\fill (-0.7246933,-2.7045923) circle (1.2pt);
\node[below right,fill=none,inner sep=2pt] at (-0.7246933,-2.7045923) {$C$};
\fill (-0.5994248,-2.2370840) circle (1.2pt);
\node[below,fill=none,inner sep=2pt] at (-0.5994248,-2.2370840) {};
\fill (0.5994248,-2.2370840) circle (1.2pt);
\node[below,fill=none,inner sep=2pt] at (0.5994248,-2.2370840) {};
\node[left,fill=none,inner sep=2pt] at (-1.6500000,-1.8000000) {$P$};
\draw[gray,dashed,thin] (-1.6000000,-1.8000000)--(-0.5994248,-2.2370840);
\node[above right,fill=none,inner sep=2pt] at (0.7794248,-1.9370840) {$Q$};
\draw[DeepBlue,thick,->] (-1.4974340,-2.4975371)--(-0.5994248,-2.2370840);
\draw[DeepBlue,thick,->] (-0.5994248,-2.2370840)--(0.5994248,-2.2370840);
\draw[DeepBlue,thick,->] (0.5994248,-2.2370840)--(1.4974340,-2.4975371);
\node[left,fill=none,inner sep=2pt] at (-1.6500000,-2.8500000) {$h$};
\draw[gray,dashed,thin] (-1.6000000,-2.8500000)--(-1.1110637,-2.6010647);
\node[left,fill=none,inner sep=2pt] at (-1.0000000,-1.2487685) {$l$};
\draw[thick] plot coordinates {(-0.1682324,-0.6278518) (-0.1649427,-0.6287240) (-0.1616484,-0.6295791) (-0.1583498,-0.6304168) (-0.1550467,-0.6312373) (-0.1517395,-0.6320404) (-0.1484281,-0.6328263) (-0.1451126,-0.6335948) (-0.1417931,-0.6343459) (-0.1384698,-0.6350796) (-0.1351426,-0.6357959) (-0.1318117,-0.6364948) (-0.1284773,-0.6371763) (-0.1251393,-0.6378402) (-0.1217979,-0.6384867) (-0.1184531,-0.6391157) (-0.1151051,-0.6397271) (-0.1117539,-0.6403211) (-0.1083997,-0.6408974) (-0.1050425,-0.6414562) (-0.1016824,-0.6419974) (-0.0983195,-0.6425210) (-0.0949540,-0.6430270) (-0.0915858,-0.6435154) (-0.0882151,-0.6439861) (-0.0848420,-0.6444392) (-0.0814666,-0.6448746) (-0.0780889,-0.6452923) (-0.0747091,-0.6456923) (-0.0713273,-0.6460746) (-0.0679435,-0.6464392) (-0.0645578,-0.6467861) (-0.0611704,-0.6471153) (-0.0577813,-0.6474267) (-0.0543906,-0.6477204) (-0.0509984,-0.6479963) (-0.0476048,-0.6482544) (-0.0442099,-0.6484948) (-0.0408138,-0.6487174) (-0.0374166,-0.6489222) (-0.0340184,-0.6491092) (-0.0306192,-0.6492784) (-0.0272192,-0.6494298) (-0.0238184,-0.6495635) (-0.0204170,-0.6496793) (-0.0170150,-0.6497773) (-0.0136126,-0.6498574) (-0.0102098,-0.6499198) (-0.0068067,-0.6499644) (-0.0034034,-0.6499911) (0.0000000,-0.6500000) (0.0034034,-0.6499911) (0.0068067,-0.6499644) (0.0102098,-0.6499198) (0.0136126,-0.6498574) (0.0170150,-0.6497773) (0.0204170,-0.6496793) (0.0238184,-0.6495635) (0.0272192,-0.6494298) (0.0306192,-0.6492784) (0.0340184,-0.6491092) (0.0374166,-0.6489222) (0.0408138,-0.6487174) (0.0442099,-0.6484948) (0.0476048,-0.6482544) (0.0509984,-0.6479963) (0.0543906,-0.6477204) (0.0577813,-0.6474267) (0.0611704,-0.6471153) (0.0645578,-0.6467861) (0.0679435,-0.6464392) (0.0713273,-0.6460746) (0.0747091,-0.6456923) (0.0780889,-0.6452923) (0.0814666,-0.6448746) (0.0848420,-0.6444392) (0.0882151,-0.6439861) (0.0915858,-0.6435154) (0.0949540,-0.6430270) (0.0983195,-0.6425210) (0.1016824,-0.6419974) (0.1050425,-0.6414562) (0.1083997,-0.6408974) (0.1117539,-0.6403211) (0.1151051,-0.6397271) (0.1184531,-0.6391157) (0.1217979,-0.6384867) (0.1251393,-0.6378402) (0.1284773,-0.6371763) (0.1318117,-0.6364948) (0.1351426,-0.6357959) (0.1384698,-0.6350796) (0.1417931,-0.6343459) (0.1451126,-0.6335948) (0.1484281,-0.6328263) (0.1517395,-0.6320404) (0.1550467,-0.6312373) (0.1583498,-0.6304168) (0.1616484,-0.6295791) (0.1649427,-0.6287240) (0.1682324,-0.6278518)};
\draw[gray,dashed,thin] (0.7000000,-0.8000000)--(0.0500000,-0.5500000);
\node[right,fill=none,inner sep=2pt] at (0.7500000,-0.8500000) {$\alpha$};
\node[below,fill=none,inner sep=2pt] at (0.0000000,-3.8000000) {湖湾内划船，岸外步行};
\end{tikzpicture}
\figcap{湖湾顶角为完整的 $\alpha$，$A,B$ 在陆地；$P,Q$ 为步行与划船的转换点。}\label{ch27:fig:s42-19-15}
\end{minipage}\end{center}
}

\newcommand{\ParallelPlateDiagram}{%
% Independent panel: original physical paths and local coordinates retained.
\begin{center}\begin{minipage}{\linewidth}\centering
\begin{tikzpicture}[>=Latex,font=\small,every node/.style={fill=none}]
\draw[thick] (-0.5000000,-0.6000000)--(-0.5000000,2.0000000);
\draw[thick] (1.0000000,-0.6000000)--(1.0000000,2.0000000);
\draw[gray,thin,->] (-2.7000000,0.0000000)--(3.0000000,0.0000000);
\fill (-2.5000000,0.0000000) circle (1.2pt);
\node[below,fill=none,inner sep=2pt] at (-2.5000000,0.0000000) {$O$};
\fill (-1.9809036,0.0000000) circle (1.2pt);
\node[below,fill=none,inner sep=2pt] at (-1.9809036,0.0000000) {$I$};
\draw[DeepBlue,thick,->] (-2.5000000,0.0000000)--(-0.5000000,0.5320431);
\draw[DeepBlue,thick,->] (-0.5000000,0.5320431)--(1.0000000,0.7929846);
\draw[DeepBlue,thick,->] (1.0000000,0.7929846)--(2.8000000,1.2718234);
\draw[gray,dashed,thin] (-1.9809036,0.0000000)--(1.0000000,0.7929846);
\node[above,fill=none,inner sep=2pt] at (2.8000000,1.2718234) {眼};
\node[above,fill=none,inner sep=2pt] at (0.2500000,1.6000000) {$n$};
\draw[gray,<->] (-2.5000000,-0.6500000)--(-0.5000000,-0.6500000);
\node[below,fill=none,inner sep=2pt] at (-1.5000000,-0.6500000) {$S$};
\draw[gray,<->] (-0.5000000,-0.6500000)--(1.0000000,-0.6500000);
\node[below,fill=none,inner sep=2pt] at (0.2500000,-0.6500000) {$D$};
\draw[gray,<->] (1.0000000,-0.6500000)--(2.8000000,-0.6500000);
\node[below,fill=none,inner sep=2pt] at (1.9000000,-0.6500000) {$L$};
\end{tikzpicture}
\figcap{玻璃板的入出射方向平行，视像 $I$ 在末段反向延长线上；$S,D,L$ 是三段轴向长度。}\label{ch27:fig:s42-19-16}
\end{minipage}\end{center}
}

\newcommand{\RainbowDiagram}{%
\begin{center}\begin{minipage}{\linewidth}\centering
\begin{tikzpicture}[>=Latex,font=\small,every node/.style={fill=none}]
\begin{scope}[shift={(-3.0000000,0.0000000)}]
\draw[thick] (0.0000000,0.0000000) circle (2);
\draw[DeepBlue,thick,->] (-2.1183502,1.7213259)--(-1.0183502,1.7213259);
\draw[gray,dashed,thin] (0.0000000,0.0000000)--(-1.0183502,1.7213259);
\fill (-1.0183502,1.7213259) circle (1.2pt);
\node[below,fill=none,inner sep=2pt] at (-1.0183502,1.7213259) {};
\node[above right,fill=none,inner sep=2pt] at (-0.4183502,2.3713259) {$A$};
\draw[gray,dashed,thin] (-1.0183502,1.7213259)--(-0.4683502,2.3213259);
\draw[gray,dashed,thin] (-1.0183502,1.7213259)--(-1.3493140,2.2807569);
\draw[thick] plot coordinates {(-1.1812862,1.9967381) (-1.1841322,1.9950344) (-1.1869604,1.9933012) (-1.1897705,1.9915389) (-1.1925622,1.9897475) (-1.1953352,1.9879273) (-1.1980892,1.9860784) (-1.2008238,1.9842011) (-1.2035388,1.9822956) (-1.2062340,1.9803620) (-1.2089089,1.9784005) (-1.2115634,1.9764115) (-1.2141971,1.9743950) (-1.2168098,1.9723514) (-1.2194011,1.9702808) (-1.2219709,1.9681834) (-1.2245187,1.9660595) (-1.2270444,1.9639093) (-1.2295477,1.9617330) (-1.2320283,1.9595310) (-1.2344860,1.9573033) (-1.2369204,1.9550502) (-1.2393313,1.9527721) (-1.2417185,1.9504691) (-1.2440817,1.9481414) (-1.2464206,1.9457894) (-1.2487351,1.9434133) (-1.2510247,1.9410133) (-1.2532894,1.9385897) (-1.2555288,1.9361428) (-1.2577428,1.9336728) (-1.2599310,1.9311799) (-1.2620933,1.9286645) (-1.2642293,1.9261269) (-1.2663390,1.9235672) (-1.2684220,1.9209858) (-1.2704782,1.9183830) (-1.2725072,1.9157589) (-1.2745090,1.9131140) (-1.2764832,1.9104485) (-1.2784297,1.9077627) (-1.2803482,1.9050568) (-1.2822386,1.9023312) (-1.2841006,1.8995861) (-1.2859341,1.8968219) (-1.2877388,1.8940388) (-1.2895146,1.8912372) (-1.2912613,1.8884173) (-1.2929786,1.8855794) (-1.2946664,1.8827240) (-1.2963245,1.8798511) (-1.2979528,1.8769613) (-1.2995510,1.8740547) (-1.3011190,1.8711317) (-1.3026567,1.8681926) (-1.3041637,1.8652377) (-1.3056401,1.8622674) (-1.3070856,1.8592819) (-1.3085001,1.8562816) (-1.3098834,1.8532668) (-1.3112353,1.8502378) (-1.3125558,1.8471950) (-1.3138447,1.8441386) (-1.3151019,1.8410691) (-1.3163271,1.8379866) (-1.3175204,1.8348917) (-1.3186815,1.8317845) (-1.3198103,1.8286655) (-1.3209067,1.8255350) (-1.3219707,1.8223932) (-1.3230020,1.8192406) (-1.3240005,1.8160775) (-1.3249663,1.8129041) (-1.3258990,1.8097210) (-1.3267988,1.8065283) (-1.3276654,1.8033265) (-1.3284987,1.8001159) (-1.3292988,1.7968968) (-1.3300654,1.7936696) (-1.3307985,1.7904346) (-1.3314981,1.7871922) (-1.3321640,1.7839428) (-1.3327962,1.7806866) (-1.3333946,1.7774240) (-1.3339592,1.7741554) (-1.3344898,1.7708811) (-1.3349865,1.7676015) (-1.3354492,1.7643169) (-1.3358777,1.7610277) (-1.3362722,1.7577342) (-1.3366325,1.7544368) (-1.3369586,1.7511359) (-1.3372505,1.7478317) (-1.3375081,1.7445247) (-1.3377315,1.7412153) (-1.3379205,1.7379036) (-1.3380751,1.7345902) (-1.3381954,1.7312754) (-1.3382814,1.7279595) (-1.3383330,1.7246429) (-1.3383502,1.7213259)};
\draw[thick] plot coordinates {(-0.7994049,1.3512409) (-0.7968135,1.3527862) (-0.7942330,1.3543498) (-0.7916636,1.3559314) (-0.7891053,1.3575309) (-0.7865583,1.3591484) (-0.7840227,1.3607838) (-0.7814987,1.3624368) (-0.7789863,1.3641076) (-0.7764857,1.3657959) (-0.7739970,1.3675018) (-0.7715204,1.3692250) (-0.7690559,1.3709656) (-0.7666036,1.3727235) (-0.7641638,1.3744985) (-0.7617365,1.3762905) (-0.7593218,1.3780996) (-0.7569199,1.3799256) (-0.7545308,1.3817684) (-0.7521547,1.3836279) (-0.7497918,1.3855040) (-0.7474420,1.3873966) (-0.7451056,1.3893057) (-0.7427826,1.3912312) (-0.7404733,1.3931729) (-0.7381776,1.3951307) (-0.7358956,1.3971046) (-0.7336276,1.3990945) (-0.7313736,1.4011003) (-0.7291338,1.4031218) (-0.7269082,1.4051589) (-0.7246969,1.4072117) (-0.7225001,1.4092799) (-0.7203178,1.4113634) (-0.7181503,1.4134623) (-0.7159975,1.4155763) (-0.7138596,1.4177053) (-0.7117367,1.4198493) (-0.7096288,1.4220081) (-0.7075362,1.4241817) (-0.7054589,1.4263699) (-0.7033970,1.4285726) (-0.7013506,1.4307897) (-0.6993198,1.4330212) (-0.6973047,1.4352668) (-0.6953054,1.4375265) (-0.6933220,1.4398002) (-0.6913547,1.4420878) (-0.6894034,1.4443891) (-0.6874683,1.4467040) (-0.6855495,1.4490324) (-0.6836471,1.4513743) (-0.6817612,1.4537295) (-0.6798918,1.4560978) (-0.6780392,1.4584792) (-0.6762032,1.4608735) (-0.6743841,1.4632806) (-0.6725820,1.4657005) (-0.6707968,1.4681329) (-0.6690288,1.4705778) (-0.6672780,1.4730351) (-0.6655445,1.4755046) (-0.6638283,1.4779861) (-0.6621296,1.4804797) (-0.6604484,1.4829851) (-0.6587849,1.4855023) (-0.6571390,1.4880310) (-0.6555110,1.4905713) (-0.6539008,1.4931229) (-0.6523085,1.4956857) (-0.6507343,1.4982597) (-0.6491781,1.5008446) (-0.6476402,1.5034404) (-0.6461205,1.5060469) (-0.6446191,1.5086640) (-0.6431361,1.5112916) (-0.6416716,1.5139296) (-0.6402256,1.5165777) (-0.6387983,1.5192359) (-0.6373896,1.5219041) (-0.6359997,1.5245821) (-0.6346286,1.5272698) (-0.6332765,1.5299670) (-0.6319432,1.5326736) (-0.6306290,1.5353896) (-0.6293339,1.5381147) (-0.6280580,1.5408488) (-0.6268012,1.5435918) (-0.6255638,1.5463435) (-0.6243456,1.5491039) (-0.6231469,1.5518727) (-0.6219676,1.5546499) (-0.6208079,1.5574353) (-0.6196677,1.5602288) (-0.6185471,1.5630302) (-0.6174463,1.5658394) (-0.6163651,1.5686562) (-0.6153038,1.5714806) (-0.6142623,1.5743123) (-0.6132407,1.5771513) (-0.6122390,1.5799974)};
\node[above,fill=none,inner sep=2pt] at (-1.6698526,2.0928693) {$i$};
\node[right,fill=none,inner sep=2pt] at (-0.3991862,1.2147334) {$r$};
\draw[DeepBlue,thick,->] (-1.0183502,1.7213259)--(1.8669753,0.7172191);
\draw[gray,dashed,thin] (0.0000000,0.0000000)--(1.8669753,0.7172191);
\draw[DeepBlue,thick,->] (1.8669753,0.7172191)--(0.3960251,-1.9603990);
\draw[gray,dashed,thin] (0.0000000,0.0000000)--(0.3960251,-1.9603990);
\draw[DeepBlue,thick,->] (0.3960251,-1.9603990)--(-0.6438926,-2.8977203);
\node[below,fill=none,inner sep=2pt] at (0.0000000,-3.3000000) {主虹：一次内反射};
\end{scope}
\begin{scope}[shift={(3.0000000,0.0000000)}]
\draw[thick] (0.0000000,0.0000000) circle (2);
\draw[DeepBlue,thick,->] (-1.7236096,1.9002924)--(-0.6236096,1.9002924);
\draw[gray,dashed,thin] (0.0000000,0.0000000)--(-0.6236096,1.9002924);
\fill (-0.6236096,1.9002924) circle (1.2pt);
\node[below,fill=none,inner sep=2pt] at (-0.6236096,1.9002924) {};
\node[above right,fill=none,inner sep=2pt] at (-0.0236096,2.5502924) {$A$};
\draw[gray,dashed,thin] (-0.6236096,1.9002924)--(-0.0736096,2.5002924);
\draw[gray,dashed,thin] (-0.6236096,1.9002924)--(-0.8262827,2.5178874);
\draw[thick] plot coordinates {(-0.7233871,2.2043392) (-0.7271910,2.2030644) (-0.7309786,2.2017420) (-0.7347494,2.2003723) (-0.7385027,2.1989554) (-0.7422379,2.1974915) (-0.7459545,2.1959809) (-0.7496518,2.1944239) (-0.7533293,2.1928206) (-0.7569865,2.1911714) (-0.7606227,2.1894764) (-0.7642373,2.1877360) (-0.7678299,2.1859504) (-0.7713998,2.1841199) (-0.7749464,2.1822448) (-0.7784693,2.1803254) (-0.7819678,2.1783619) (-0.7854414,2.1763548) (-0.7888896,2.1743042) (-0.7923118,2.1722106) (-0.7957076,2.1700743) (-0.7990762,2.1678955) (-0.8024173,2.1656747) (-0.8057303,2.1634122) (-0.8090146,2.1611083) (-0.8122698,2.1587635) (-0.8154954,2.1563780) (-0.8186908,2.1539522) (-0.8218555,2.1514866) (-0.8249891,2.1489815) (-0.8280910,2.1464374) (-0.8311608,2.1438545) (-0.8341979,2.1412333) (-0.8372020,2.1385743) (-0.8401725,2.1358778) (-0.8431089,2.1331443) (-0.8460109,2.1303742) (-0.8488778,2.1275680) (-0.8517094,2.1247260) (-0.8545052,2.1218487) (-0.8572646,2.1189366) (-0.8599873,2.1159902) (-0.8626729,2.1130098) (-0.8653209,2.1099961) (-0.8679309,2.1069493) (-0.8705025,2.1038701) (-0.8730353,2.1007589) (-0.8755289,2.0976161) (-0.8779829,2.0944424) (-0.8803969,2.0912381) (-0.8827705,2.0880038) (-0.8851034,2.0847401) (-0.8873953,2.0814473) (-0.8896456,2.0781261) (-0.8918542,2.0747769) (-0.8940205,2.0714002) (-0.8961444,2.0679967) (-0.8982255,2.0645669) (-0.9002634,2.0611112) (-0.9022577,2.0576302) (-0.9042083,2.0541245) (-0.9061148,2.0505946) (-0.9079769,2.0470411) (-0.9097943,2.0434646) (-0.9115668,2.0398655) (-0.9132939,2.0362445) (-0.9149755,2.0326021) (-0.9166114,2.0289389) (-0.9182012,2.0252555) (-0.9197446,2.0215525) (-0.9212416,2.0178304) (-0.9226917,2.0140899) (-0.9240949,2.0103314) (-0.9254508,2.0065557) (-0.9267592,2.0027632) (-0.9280201,1.9989547) (-0.9292330,1.9951306) (-0.9303980,1.9912916) (-0.9315147,1.9874384) (-0.9325830,1.9835714) (-0.9336028,1.9796913) (-0.9345738,1.9757988) (-0.9354960,1.9718944) (-0.9363691,1.9679787) (-0.9371931,1.9640524) (-0.9379678,1.9601161) (-0.9386931,1.9561704) (-0.9393689,1.9522159) (-0.9399950,1.9482532) (-0.9405714,1.9442830) (-0.9410980,1.9403059) (-0.9415747,1.9363225) (-0.9420014,1.9323334) (-0.9423781,1.9283393) (-0.9427046,1.9243408) (-0.9429811,1.9203385) (-0.9432073,1.9163331) (-0.9433833,1.9123251) (-0.9435090,1.9083152) (-0.9435844,1.9043041) (-0.9436096,1.9002924)};
\draw[thick] plot coordinates {(-0.4895335,1.4917295) (-0.4862970,1.4928059) (-0.4830691,1.4939078) (-0.4798501,1.4950354) (-0.4766401,1.4961885) (-0.4734394,1.4973669) (-0.4702481,1.4985708) (-0.4670664,1.4997999) (-0.4638946,1.5010542) (-0.4607329,1.5023336) (-0.4575814,1.5036381) (-0.4544403,1.5049675) (-0.4513099,1.5063218) (-0.4481904,1.5077008) (-0.4450818,1.5091046) (-0.4419846,1.5105330) (-0.4388987,1.5119859) (-0.4358244,1.5134633) (-0.4327620,1.5149650) (-0.4297116,1.5164909) (-0.4266734,1.5180410) (-0.4236475,1.5196151) (-0.4206343,1.5212132) (-0.4176338,1.5228351) (-0.4146463,1.5244808) (-0.4116719,1.5261501) (-0.4087109,1.5278430) (-0.4057633,1.5295593) (-0.4028295,1.5312989) (-0.3999096,1.5330618) (-0.3970038,1.5348477) (-0.3941122,1.5366566) (-0.3912350,1.5384885) (-0.3883725,1.5403430) (-0.3855248,1.5422203) (-0.3826920,1.5441200) (-0.3798744,1.5460422) (-0.3770722,1.5479867) (-0.3742854,1.5499533) (-0.3715144,1.5519419) (-0.3687592,1.5539525) (-0.3660200,1.5559849) (-0.3632971,1.5580389) (-0.3605905,1.5601145) (-0.3579005,1.5622115) (-0.3552272,1.5643297) (-0.3525708,1.5664691) (-0.3499314,1.5686295) (-0.3473092,1.5708108) (-0.3447045,1.5730128) (-0.3421173,1.5752353) (-0.3395477,1.5774784) (-0.3369961,1.5797417) (-0.3344625,1.5820252) (-0.3319471,1.5843288) (-0.3294500,1.5866522) (-0.3269715,1.5889953) (-0.3245116,1.5913580) (-0.3220706,1.5937402) (-0.3196485,1.5961417) (-0.3172455,1.5985623) (-0.3148618,1.6010019) (-0.3124975,1.6034603) (-0.3101528,1.6059374) (-0.3078279,1.6084330) (-0.3055228,1.6109469) (-0.3032377,1.6134791) (-0.3009727,1.6160293) (-0.2987281,1.6185974) (-0.2965039,1.6211833) (-0.2943003,1.6237867) (-0.2921174,1.6264074) (-0.2899554,1.6290455) (-0.2878143,1.6317006) (-0.2856944,1.6343725) (-0.2835958,1.6370613) (-0.2815185,1.6397665) (-0.2794628,1.6424882) (-0.2774287,1.6452261) (-0.2754164,1.6479800) (-0.2734260,1.6507499) (-0.2714576,1.6535354) (-0.2695114,1.6563364) (-0.2675875,1.6591528) (-0.2656860,1.6619844) (-0.2638070,1.6648309) (-0.2619506,1.6676923) (-0.2601170,1.6705683) (-0.2583063,1.6734588) (-0.2565185,1.6763635) (-0.2547538,1.6792823) (-0.2530124,1.6822150) (-0.2512942,1.6851615) (-0.2495995,1.6881215) (-0.2479284,1.6910948) (-0.2462808,1.6940813) (-0.2446570,1.6970808) (-0.2430570,1.7000930) (-0.2414810,1.7031179) (-0.2399291,1.7061551) (-0.2384012,1.7092046)};
\node[above,fill=none,inner sep=2pt] at (-1.2310181,2.3402411) {$i$};
\node[right,fill=none,inner sep=2pt] at (-0.0999029,1.2955351) {$r$};
\draw[DeepBlue,thick,->] (-0.6236096,1.9002924)--(1.8903165,0.6532255);
\draw[gray,dashed,thin] (0.0000000,0.0000000)--(1.8903165,0.6532255);
\draw[DeepBlue,thick,->] (1.8903165,0.6532255)--(0.6826820,-1.8798791);
\draw[gray,dashed,thin] (0.0000000,0.0000000)--(0.6826820,-1.8798791);
\draw[DeepBlue,thick,->] (0.6826820,-1.8798791)--(-1.8689827,-0.7119717);
\draw[gray,dashed,thin] (0.0000000,0.0000000)--(-1.8689827,-0.7119717);
\draw[DeepBlue,thick,->] (-1.8689827,-0.7119717)--(-2.7504472,0.3756957);
\node[below,fill=none,inner sep=2pt] at (0.0000000,-3.3000000) {次虹：两次内反射};
\end{scope}
\end{tikzpicture}
\figcap{主虹与次虹的球内路径；虚线为法线，$i,r$ 在入射点 $A$ 测量。离球后的箭头指向接收这束光的观察者。\label{ch27:fig:rainbow}}
\end{minipage}\end{center}
}


光学是一门古老而至今仍十分活跃的学科。现存最早的人工透镜可追溯到公元前七世纪；从以牛顿为代表的近代物理学开始，光学理论和光学仪器的制造一同进入快速发展时期。今天，光学不仅是基础物理的重要组成部分，也深刻影响着通信、精密测量、医学、天文、材料和信息技术。

贯穿光学发展史的核心问题是：光究竟是什么？牛顿（Isaac Newton）曾把光理解为高速运动的粒子流；同一时期，惠更斯（Christiaan Huygens）主张光是一种波。十九世纪初，杨（Thomas Young）的干涉实验和菲涅尔（Augustin-Jean Fresnel）的理论使波动说重新兴起；麦克斯韦（James Clerk Maxwell）随后指出光是横向电磁波，不同颜色对应不同频率或波长，而真空中的传播速度相同。二十世纪初，爱因斯坦（Albert Einstein）的光量子理论揭示了能量交换的离散性。波粒二象性并不表示单个光子会在“传播时是经典波、探测时是经典粒子”之间切换：光应由量子态或量子化电磁场统一描述，概率振幅可以叠加并产生干涉，探测器则记录一次次离散的能量交换事件。经典波和经典粒子只是不同实验条件下的近似图像；量子电动力学中，光子是电磁场的量子化激发。


光学仪器的发展构成了另一条同样重要的主线。平面镜的平整度、透镜和反射镜的面形精度、材料纯度以及装调误差都会改变仪器的性能。现实中不存在绝对完美的光学系统，设计者必须在功能、精度、口径、重量和成本之间权衡。历史人物与代表性器件提示我们：光的理论与实验仪器始终相互推动。

\historyfigure{optics-history-built-in-v1.png}{光学史人物与代表性实验器件：粒子说、波动说、棱镜色散、望远镜与双缝实验。}

当光学元件和传播尺度远大于光的波长时，波动光学可以取极短波长极限，得到更简单的\textbf{几何光学}。几何光学用“光线”表示能量传播方向，但光线不是有实体的细线，也不是经典粒子的真实轨迹；在各向同性介质中，它是波面法线在空间中的积分曲线。

因此几何光学有明确的适用边界：器件尺寸、曲率半径和传播中被分辨的细节都应远大于波长。接近小孔、尖锐边缘或亚波长结构时，单靠光线会漏掉衍射；两路光重新相遇时，单靠路径也不能判断明暗。后两讲将处理这些波动效应，本讲则先在短波长极限内建立可用于追踪方向的规律。

在任一时刻，把光场中相位相同的点连接起来，就得到等相位面（波面）。在各向同性透明介质中，光能沿波面的法线方向传播。点光源在均匀各向同性介质中的波面是一组同心球面，所以光线是从光源出发的直线。这给出几何光学的第一条基本规律：\textbf{光在均匀介质中沿直线传播}。

若介质不均匀，波速随位置改变，波面各部分前进速度不同，光线便会弯曲。图~\ref{ch27:fig:s42-19-1} 中，大气密度通常向地面方向增大，折射率也相应增大；接近地平线的太阳光逐层向法线方向偏折。观察者沿最后一段光线反向延长，会把太阳判断在真实位置之上，因此太阳已经略低于几何地平线时仍可能被看见。

\AtmosphereWavefrontPanel



\subsection{从波长、折射率到光程}
光的频率记为 $\nu$，角频率为 $\omega=2\pi\nu$。真空波长 $\lambda_0=c/\nu$ 与介质内波长不是同一个量。在静止界面上，入射波、反射波和折射波必须以相同频率振动，否则界面上的电场边界条件不能在每个时刻同时满足。若透明介质的相折射率为 $n$，则相速度 $v_p=c/n$，介质内波长为
\[
 \lambda=\frac{v_p}{\nu}=\frac{\lambda_0}{n}.
\]
所以“进入玻璃后波长变短”不表示光的频率降低。色散指 $n$ 随频率而变；比较不同颜色时，要分别使用各自的折射率。

一段长度为 $s$ 的均匀介质，使单色波积累的传播相位为
\[
 ks=\frac{2\pi}{\lambda}s=\frac{2\pi}{\lambda_0}ns.
\]
因此把 $ns$ 称为这段路的光程：它把“走过多长”和“介质使相位变化多快”合在一起。依次通过若干介质时，光程相加为 $\sum_jn_js_j$；折射率连续变化时才把求和改为积分。光程的单位仍是米，折射率没有量纲。它不是把折射后的折线路径拉直所得到的几何长度。

本章的普通玻璃、水和空气先按各向同性介质处理。在这种介质中，波面法线、波矢与能量传播方向一致；晶体中的非寻常光则不一定如此，第30章再区分波法线与能流方向。脉冲通过明显色散的介质时，能量包络的传播时间涉及群速度，也不能直接把相折射率代入 $t=ns/c$。


\section{反射定律与镜面光路}



光照到平整界面时，一部分能量返回原介质，称为反射。在入射点作界面的法线；入射光线与法线的夹角称为入射角 $i$，反射光线与法线的夹角称为反射角 $r$。入射光线、法线和反射光线共面，并且
\begin{equation}
  i=r.
\end{equation}
角度必须相对\textbf{法线}测量。对曲面，应先在入射点作切平面，再以切平面的垂线作为法线。

两块平面镜夹角为 $\alpha$ 时，连续两次反射后，出射方向相对入射方向的总转角为
\begin{equation}
  \Delta\varphi=2\alpha,
\end{equation}
其大小与第一束光的入射角无关。$\alpha=0$ 对应平行镜，可用于平移光路；$\alpha=\pi/2$ 时出射方向与入射方向相反，构成直角反射器。月球激光测距所用的角锥反射器正是这一思想的三维推广。

这个结论可以直接用方向角核验。设一条光线相对固定轴的有向角为 $\varphi$，镜面方向角为 $\beta$。镜面反射把光线关于镜面对称，故一次反射后
\[
 \varphi'=2\beta-\varphi.
\]
若两镜方向角依次为 $\beta_1,\beta_2$，两次反射后
\[
 \varphi''=2\beta_2-(2\beta_1-\varphi)
 =\varphi+2(\beta_2-\beta_1).
\]
因此方向转角确为 $2\alpha$，而不是 $\pi-2\alpha$。后一表达式有时表示两条无向直线间取补角后的夹角，不能当作光的有向传播方向变化。图~\ref{ch27:fig:s42-19-3} 中入射箭头、两次反射点和出射箭头正是按这一关系绘制的。

\historyfigure[0.96\linewidth]{reflection-applications-built-in-v1.png}{反射定律实验；潜望镜与月面激光反射器。\label{ch27:fig:s42-19-2}}
图~\ref{ch27:fig:s42-19-2} 将反射定律放到实验和装置中：潜望镜利用多次反射改变光路，月面反射器用于激光往返测距；逐个反射面都要以该面的法线量角。

\subsection{先定法线，再追踪反射方向}
入射角从法线量起，不从镜面量起。例如光线与镜面成 $20^\circ$，入射角就是 $70^\circ$。角度数值正确还不够：反射光必须留在入射侧，并沿离开反射点的方向传播。把一条直线的两个相反方向混为一谈，往往会得到“线平行但箭头错误”的光路。

设入射传播方向的单位矢量为 $\vect s$，镜面单位法向为 $\vect n$。分解成沿面与垂直面两部分：
\[
 \vect s=\bigl[\vect s-(\vect s\cdot\vect n)\vect n\bigr]
       +(\vect s\cdot\vect n)\vect n.
\]
反射后沿面分量不变，法向分量反号，因而
\begin{equation}\label{ch27:eq:depth-reflection-vector}
 \vect s_{\rm r}=\vect s-2(\vect s\cdot\vect n)\vect n.
\end{equation}
换用 $-\vect n$ 不改变结果。这一式既可检查平面镜，也可检查曲面镜，只需把 $\vect n$ 换成入射点处的局部法线。球面法线必须经过球心；它一般不是图上的竖直线或水平线。

连续两次反射时，应把第一次反射后的传播方向作为第二次的入射方向。用镜像法展开折线只是帮助计算：镜后虚像和虚线延长线并不表示光实际穿进镜后。求光斑运动时还要先确定最后一条反射线与屏的交点，再对交点位置求时间导数，不能把镜子的角速度直接当作光斑的线速度。


\ReflectionPanel


\SphericalMirrorRotatingDiagram
\Needspace{7\baselineskip}\begin{origexamplebox}{【球面镜】}\label{ch27:ex:s42-19-1}
有一个半径为 $R$ 的球形反射面，圆心位于 $C$，$CV$ 为其对称轴。一束与对称轴平行、距轴为 $d$ 的入射光经球面反射后与对称轴交于 $M$。求 $M$ 与球面顶点 $V$ 之间的距离 $s$。
\end{origexamplebox}\par\nopagebreak[3]
\begin{xesolution}
取子午截面。设入射点为 $P$，令 $u=\sqrt{R^2-d^2}$。从顶点 $V$ 向球心方向量，$P$ 的轴向坐标为 $R-u$。半径 $CP$ 是法线，记它与轴的夹角为 $\alpha$，则 $\sin\alpha=d/R$、$\cos\alpha=u/R$。入射光平行于轴，故反射光与轴的夹角为 $2\alpha$。由图~\ref{ch27:fig:s42-19-11} 左图中的直角三角形，
\[
 s=(R-u)+\frac{d}{\tan 2\alpha}.
\]
又 $\tan2\alpha=2du/(u^2-d^2)$，整理得
\[
 \boxed{s=R-\frac{R^2}{2\sqrt{R^2-d^2}}}.
\]
上式的 $s$ 是从 $V$ 向 $C$ 为正的有向坐标；题目所问几何距离为 $|s|$，参数域为 $|d|<R$。当 $|d|<\sqrt3R/2$ 时 $s>0$，超过此高度时轴交点移到顶点另一侧，$s<0$。当 $d\to0$ 时 $s\to R/2$；不同高度给出不同交点，说明球面镜存在球差。$d=0$ 的轴上光与轴重合，焦点位置须由近轴极限定义。
\end{xesolution}

\Needspace{7\baselineskip}\begin{origexamplebox}{【转动的镜子】}\label{ch27:ex:s42-19-2}
一竖直墙上 $A$ 点发出水平光束，平面镜 $OB$ 绕过 $O$ 的水平轴以角速度 $\omega$ 转动。光束经镜面反射后回到墙上形成光点 $D$。已知此刻镜面与墙的夹角为 $\theta$，$AO=d$，忽略光的传播时间，求光点在墙上的移动速度。
\end{origexamplebox}\par\nopagebreak[3]
\begin{xesolution}
%<optics-review:mirror-domain>
先限定能够在墙上看到反射光点的姿态。取图示镜面位于墙左侧、$0<\theta<\pi/2$，入射光从墙向左传播；以下把 $\omega=\dot\theta$ 作为带符号角速度。若题给的是角速度大小，须先由转向确定正负。反射方向相对向右的水平线转过 $2\theta$，其单位方向矢量为 $(\cos2\theta,\sin2\theta)$。因此光要从左侧镜面返回右侧墙面，必须有
\[
 \cos2\theta>0,\qquad 0<\theta<\frac\pi4.
\]
在 $\theta=\pi/4$ 时反射光平行于墙；超过这个角度，反射光向左离开，直线延长线与墙的交点不能再解释成实际光点。这里还默认入射点位于镜面的有效长度内。

%</optics-review:mirror-domain>

在垂直转轴的截面内令墙为 $x=0$，$O=(0,0)$，$A=(0,d)$。镜面直线可写为 $x=-y\tan\theta$，水平入射线 $y=d$ 与镜面交于 $C=(-d\tan\theta,d)$。镜面转过 $\theta$ 时反射方向转过 $2\theta$，故从 $C$ 到墙的竖直位移为 $d\tan\theta\tan2\theta$。于是
\[
 y_D=d+d\tan\theta\tan2\theta.
\]

%<optics-review:mirror-simplify>
这里 $d\tan\theta$ 是反射点到墙的水平距离，乘上反射方向的斜率 $\tan2\theta$ 才是竖直位移。利用
\[
 1+\tan\theta\tan2\theta
 =\frac{\cos\theta\cos2\theta+\sin\theta\sin2\theta}
 {\cos\theta\cos2\theta}
 =\frac1{\cos2\theta},
\]
也可把位置写成 $y_D=d\sec2\theta$，从而直接得到
\[
 v_D=2d\omega\sec2\theta\tan2\theta.
\]
这是下式的等价形式。在上述允许区间内，$\omega>0$ 时光点上移，$\omega<0$ 时下移；其量纲为长度乘角速度，即速度。

%</optics-review:mirror-simplify>
利用 $\dot\theta=\omega$ 求导，得到沿墙的瞬时速度
\[
 \boxed{v_D=d\omega\!\left(\sec^2\theta\tan2\theta+2\tan\theta\sec^22\theta\right)}.
\]
式子的正负表示移动方向；题目若只问速率，应取绝对值。$\theta$ 接近使反射光平行于墙的位置时，$\tan2\theta$ 发散，光点速度可很大，这并不表示物质或信息以该速度传播。
\end{xesolution}

\section{折射定律与多界面光路}



光从一种透明介质进入另一种透明介质时，传播方向一般会改变。折射光线仍位于入射平面内。若界面两侧折射率分别为 $n_1,n_2$，入射角和折射角分别为 $\theta_1,\theta_2$，则斯涅尔定律（Snell's law）为
\begin{equation}
  n_1\sin\theta_1=n_2\sin\theta_2.
\end{equation}
折射率定义为 $n=c/v$，其中 $c$ 为真空光速，$v$ 为介质中的相速度。在本讲讨论的普通透明介质中 $n>1$。同一介质的折射率通常随波长变化，所以复色光会发生色散。近轴条件 $|\theta|\ll 1$（角度用弧度表示）下，$\sin\theta\simeq\theta$，斯涅尔定律近似为
\begin{equation}
  n_1\theta_1\simeq n_2\theta_2.
\end{equation}



\subsection{多次折射中什么可以消去}
折射角总是相对于当前界面的法线定义。若各界面互相平行，法线方向相同，一段介质内的出射角就等于下一界面的入射角，于是
\[
 n_1\sin i_1=n_2\sin i_2=\cdots=n_k\sin i_k.
\]
中间折射率可以从“最终方向”的关系中消去，但中间层的厚度、侧向位移和光程不能消去。若两界面不平行，例如棱镜的两个侧面，必须先用几何关系换算第二次入射角，再用折射定律。

以厚度为 $t$ 的平行玻璃板为例，板外同为空气，空气侧角为 $i$，板内角为 $r$。第二次折射满足 $n\sin r=\sin i'$，因此 $i'=i$：出射光与入射光同向平行。板内光线的几何长度为 $t/\cos r$，相对原直线的垂直侧移量则为
\[
 b=\frac{t\sin(i-r)}{\cos r}.
\]
这里 $t\tan i-t\tan r$ 是在平行于板面的方向上比较得到的交点差，再乘 $\cos i$ 才得到两条平行光线的最短距离 $b$。当 $n\to1$ 或 $t\to0$ 时，$b\to0$；正入射时也没有侧移，但仍会产生光程变化。

处理球壳或圆柱时，先连“球心或圆心—入射点”得到法线。光线到中心的垂距、光线与半径的夹角和球半径之间满足直角三角形关系；先建立这一关系，才能把折射定律与“碰到内壁”“与内球相切”等条件连接起来。


\CylinderDiagram

\Needspace{7\baselineskip}\begin{origexamplebox}{【圆柱折射】}\label{ch27:ex:s42-19-3}
一个半径 $R=2\,\mathrm m$ 的透明圆柱体，其一侧是反射面。一束光从左侧空气入射，经折射和反射后，出射光与入射光平行；两平行光线的距离 $d=2\,\mathrm m$。求透明介质的折射率。
\end{origexamplebox}\par\nopagebreak[3]
\begin{xesolution}
在垂直圆柱轴的截面内分析。图~\ref{ch27:fig:s42-19-12} 中两外部光线关于过圆心 $C$ 的水平轴对称，入射点 $A$ 的高度为 $d/2=R/2$。取 $C=(0,0)$，便有 $A=(-\sqrt3R/2,R/2)$、反射点 $P=(R,0)$，因此相对向内法线 $AC$ 的入射角 $i=30^\circ$。内部线段 $AP$ 相对水平轴向下的倾角为 $\beta$，其中
\[
 \tan\beta=\frac{R/2}{R+\sqrt3R/2}=2-\sqrt3=\tan15^\circ.
\]
所以 $\beta=15^\circ$，折射角为法线与内部光线的夹角 $r=i-\beta=15^\circ$。在 $P$ 处按水平法线反射后，光线到达对称点 $B$，再折射为水平向左的出射光，方向也与题设一致。斯涅尔定律给出
\[
 1\cdot\sin30^\circ=n\sin15^\circ,\qquad
 \boxed{n=\frac{\sin30^\circ}{\sin15^\circ}=\frac{\sqrt6+\sqrt2}{2}\approx1.93}.
\]
图~\ref{ch27:fig:s42-19-12} 同时画出了两次折射点和右侧反射点；若只凭“平行”而忽略圆柱半径法线的位置关系，角度会误判。
\end{xesolution}

% Relocated AbsorbingShellDiagram to its own solution.

\Needspace{7\baselineskip}\begin{origexamplebox}{【大小球】}\label{ch27:ex:s42-19-4}
用折射率为 $n$ 的透明物质制成内、外半径分别为 $a,b$ 的空心球，内表面涂有能完全吸收光的物质。平行光射向此球时，球吸收的入射光束横截面积多大？
\end{origexamplebox}\par\nopagebreak[3]
\begin{xesolution}
\AbsorbingShellDiagram

先采用几何光学模型，忽略外表面的菲涅尔反射损耗，求能到达吸收层的入射光束面积，而不是包含反射损耗的功率吸收截面。取图~\ref{ch27:fig:s42-19-13} 中恰与内球在 $T$ 相切的临界光线。设它进入外球面后的折射角为 $r$，则直线到球心的距离为 $b\sin r=a$。设这条光线在入射平面内的入射高度（入射线到球心的垂直距离）为 $\rho$，空气侧有 $\rho=b\sin i$。由斯涅尔定律 $\sin i=n\sin r$，所以
\[
 \rho=b\sin i=nb\sin r=na.
\]
入射高度小于 $\rho$ 的光线会碰到吸收层，故
\[
 \boxed{S_{\rm abs}=\pi\min\{b^2,n^2a^2\}}.
\]
若 $na\ge b$，所有射入外球几何截面的光都能到达内层，面积饱和为 $\pi b^2$；这一截断条件不能省略。
\end{xesolution}

\Needspace{7\baselineskip}\begin{origexamplebox}{【平板玻璃砖】}\label{ch27:ex:s42-19-9}
观察者距折射率为 $n$、厚度为 $D$ 的平行玻璃板近表面为 $L$。发光点位于玻璃板另一侧，距远表面为 $S$。求观察者看到的发散光线反向延长后的视距离。
\end{origexamplebox}\par\nopagebreak[3]
\begin{xesolution}
\ParallelPlateDiagram

%<optics-review:plate-reference>
参看图~\ref{ch27:fig:s42-19-16}：$O$ 是光源，$I$ 在玻璃出射光的反向延长线上，眼睛在玻璃另一侧。$S,D,L$ 分别是光源到玻璃、玻璃内部、玻璃到眼睛的轴向长度。应从眼睛收到的末段方向反推位置，不是把各段的相位光程 $n\ell$ 相加。

%</optics-review:plate-reference>

取一条近轴光线。在空气中角度为 $\theta$，玻璃中角度为 $\theta_g$。由近轴斯涅尔定律 $\theta_g\simeq\theta/n$。光线在三个区段产生的横向位移为
\[
 y\simeq L\theta+D\theta_g+S\theta
  =\theta\left(L+\frac Dn+S\right).
\]
观察者按空气中的末段方向反向延长，虚像点到观察者的距离就是 $y/\theta$：
\[
 \boxed{L_{\rm view}=L+\frac Dn+S}.
\]
所以玻璃的视厚度为 $D/n$，比真实厚度小。该结果依赖近轴条件；大角度光线的反向延长线不再严格交于同一点。
\end{xesolution}

\section{棱镜偏向、色散与彩虹}

对顶角为 $A$、置于空气中的三棱镜，记两面外侧入射/出射角为 $i_1,i_2$，棱镜内的折射角为 $r_1,r_2$。由几何关系
\begin{equation}
 r_1+r_2=A,\qquad \delta=i_1+i_2-A.
\end{equation}
在对称光路 $i_1=i_2$、$r_1=r_2=A/2$ 处偏向角取极小值。普通玻璃中蓝光折射率较大，因而偏折通常比红光更强。牛顿的棱镜实验正是由此说明白光由多种颜色组成。

棱镜计算必须把两个界面依次写全。空气折射率取 $1$ 时，第一面与第二面分别满足
\[
 \sin i_1=n\sin r_1,\qquad n\sin r_2=\sin i_2.
\]
给定 $A,n$ 和 $i_1$ 后，先由第一式求 $r_1$，再用 $r_2=A-r_1$ 求 $i_2$，最后才代入 $\delta=i_1+i_2-A$。图示通常取两内部角均为正；若光线越过法线，应使用有向角。令 $r_c=\arcsin(1/n)$，从指定两面贯穿的必要角度条件是 $|r_1|,|r_2|\le r_c$，同时要确认中间光段没有先碰到底面。若第二面上 $|r_2|>r_c$，便发生全反射。等号表示掠射极限，不能把它当作具有有限法向能流的出射光。

把入、出方向交换，光路可逆，因此交换 $r_1,r_2$ 不改变偏向角；但这种对称性本身并不能证明极小。令 $f(r)=\arcsin(n\sin r)$，则 $\delta=f(r_1)+f(A-r_1)-A$。函数 $f'(r)=n\cos r/\sqrt{1-n^2\sin^2r}$ 为偶函数，并随 $|r|$ 增大而增大。由于 $A>0$，$r_1<A/2$ 时有 $|r_1|<|A-r_1|$，故 $\delta'<0$；$r_1>A/2$ 时相反。于是当 $A<2r_c$、对称光路可贯穿棱镜时，唯一极小值在 $r_1=r_2=A/2$，即
\[
 \delta_{\min}=2\arcsin\!\left(n\sin\frac A2\right)-A.
\]
实验中缓慢转动棱镜时，光斑先向一个方向移动、到达最小偏向位置后再反向；这比仅凭一张静态光路图更容易辨认极小值。

\RefractionPrismPanel

图\ref{ch27:fig:deviation}给出固定折射率下的偏向角曲线。入射角增大时偏向角先减小后增大，最低点对应内部对称光路；它与色散图回答的是两个不同问题。
\subsection{最小偏向与色散测量的条件}
棱镜公式中的 $A$ 是两个工作面的夹角，不一定等于图上最醒目的角。设内部两次折射角分别为 $r_1,r_2$，有 $r_1+r_2=A$。在可逆、对称的最小偏向光路中，$r_1=r_2=A/2$、$i_1=i_2=i$，故
\[
 \delta_{\min}=2i-A,\qquad
 n=\frac{\sin[(A+\delta_{\min})/2]}{\sin(A/2)}
\]
（外界为空气）。实验上固定棱镜并只改变波长时，不同颜色一般不同时满足最小偏向条件；不能把同一个入射角问题误做成“每种颜色都重新调整到最小偏向”。

薄棱镜近似还要求 $A,i,r$ 都足够小。把正弦近似成角度后，$i_1\simeq nr_1$、$i_2\simeq nr_2$，于是
\[
 \delta=i_1+i_2-A\simeq(n-1)A,
 \qquad \delta_{\rm violet}-\delta_{\rm red}
 \simeq(n_{\rm violet}-n_{\rm red})A.
\]
这说明色散不仅来自“有棱镜”，更来自不同频率的折射率不同。若需要较大的偏向角或高精度结果，必须回到两次折射的精确关系，不能继续用线性近似。


\PrismDeviationDiagram



\Needspace{7\baselineskip}\begin{origexamplebox}{【最小偏向角】}\label{ch27:ex:s42-19-11}
一三棱镜顶角为 $A$、折射率为 $n$，光由一个面入射并从另一面射出。求偏向角的可能范围。
\end{origexamplebox}\par\nopagebreak[3]
\begin{xesolution}
取空气中 $n>1$、$0<A<\pi$，讨论从指定两侧面贯穿且未被底面截断的光路。为覆盖小顶角情形，内部角 $r_1,r_2$ 采用有向角，满足 $r_1+r_2=A$；负值表示光线越过对应法线。外侧角为 $i_j=\arcsin(n\sin r_j)$，偏向角 $\delta=i_1+i_2-A$。
令 $r_c=\arcsin(1/n)$，两面透射要求 $|r_j|\le r_c$，于是
\[
 A-r_c\le r_1\le r_c.
\]
若 $A>2r_c$，区间为空，没有贯穿两面的透射光路；若 $A=2r_c$，只剩双掠射极限。以下取 $0<A<2r_c$。
写 $F(r)=\arcsin(n\sin r)$，其导数
\[
 F'(r)=\frac{n\cos r}{\sqrt{1-n^2\sin^2r}}
\]
为偶函数，随 $|r|$ 严格增加（$n>1$）。因此
$\delta'(r_1)=F'(r_1)-F'(A-r_1)$ 在 $r_1<A/2$ 时为负，在 $r_1>A/2$ 时为正。
最小值确实在 $r_1=r_2=A/2$，而不只由可逆性猜出：
\[
 \boxed{\delta_{\min}=2\arcsin\!\left(n\sin\frac A2\right)-A}.
\]
偏向角上界出现在某一外侧光线掠射界面。取 $r_1=r_c$，$r_2=A-r_c$，得到
\[
 \boxed{\delta_{\max}=\frac\pi2+\arcsin[n\sin(A-r_c)]-A}.
\]
这一式在 $A<r_c$ 时仍成立，此时第二个反正弦为负，不能擅自截为零。实际非掠射光路的范围为 $\delta_{\min}\le\delta<\delta_{\max}$；若允许极限光线，则右端也可写等号。特定入射位置、有限底面和光阑还会缩小可达范围，以上是可选择入射位置的两侧面几何范围。角度统一用弧度。
\end{xesolution}

\RainbowDiagram

\Needspace{7\baselineskip}\begin{origexamplebox}{【彩虹成因】}\label{ch27:ex:s42-19-12}
雨后出现彩虹。把空中的水滴看作折射率为 $n$ 的球，解释主虹的形成，并求其特征方向。
\end{origexamplebox}\par\nopagebreak[3]
\begin{xesolution}
主虹对应“入射折射—一次内反射—出射折射”的光路，如图\ref{ch27:fig:rainbow}左侧路径。入射点 $A$ 的半径是法线。设空气侧入射角为 $i$、水滴内折射角为 $r$，$\sin i=n\sin r$。光线相对入射方向的总偏折为
\[
 D_1=\pi+2i-4r.
\]

%<optics-review:rainbow-turns>
这个偏折角可逐段相加：两次折射各使方向转过 $i-r$，球内一次反射使方向转过 $\pi-2r$，所以 $D_1=2(i-r)+(\pi-2r)$。图中不同入射高度对应不同 $i$，并非每条光都偏折相同角度。

%</optics-review:rainbow-turns>
大量邻近光线在 $\dd D_1/\dd i=0$ 的方向聚集，形成明亮的虹带。由 $\cos i=n\cos r\,\dd r/\dd i$ 得
\[
 2-4\frac{\cos i}{n\cos r}=0
 \quad\Longrightarrow\quad
 \sin^2i=\frac{4-n^2}{3}.
\]

%<optics-review:rainbow-angles>
上面的代数步骤是：驻值条件给出 $2\cos i=n\cos r$，两边平方，再用 $n\sin r=\sin i$，便有
\[
 4(1-\sin^2i)=n^2-\sin^2i,
 \qquad 3\sin^2i=4-n^2.
\]
普通水滴满足 $1<n<2$，故驻值在非掠射的允许角域内。它对应最小偏折：在 $i=0$ 附近 $D_1'=2-4/n<0$，接近掠射时 $D_1'\to2>0$，而驻值方程在 $0<i<\pi/2$ 仅有一个解。

观察到的虹半径不是 $D_1$，而是视线与反太阳方向的夹角
\[
 \rho=\pi-D_1=4r-2i.
\]
取 $n=4/3$，依次求得 $i\approx59.39^\circ$、$r\approx40.20^\circ$，于是 $D_1\approx137.97^\circ$、$\rho\approx42.03^\circ$。驻值附近入射角改变较多而出射方向改变很少，光线因此集中；实际虹的亮度与宽度还受太阳有限角宽、水滴大小及波动效应影响，并非无限亮的一条线。

%</optics-review:rainbow-angles>
对水的可见光折射率代入，主虹的观察半径约为 $42^\circ$。折射率随波长变化，红光偏折较小，故主虹红色在外、紫色在内。单个水滴只把某一颜色送入某一方向；整条虹来自视线方向不同的大量水滴。
\end{xesolution}

\Needspace{7\baselineskip}\begin{origexamplebox}{【广义彩虹】}\label{ch27:ex:s42-19-13}
网页和自然界中还能看到多种类似彩虹的彩色弧带或光环。怎样判断并解释这些图样的成因？
\end{origexamplebox}\par\nopagebreak[3]
\begin{xesolution}
不能把所有彩色圆弧都归结为同一种几何光路，应先看光源方向、弧心位置、颜色次序和介质形态。
\begin{itemize}
  \item 次虹由水滴内两次反射形成，偏折 $D_2=2\pi+2i-6r$；它位于主虹外侧、亮度较弱，颜色次序与主虹相反，主虹与次虹之间形成较暗的 Alexander 带。
  \item 雾虹仍与水滴有关，但液滴很小，衍射展宽压过色散，因而弧带宽而近白色。
  \item 日晕和月晕主要由高空六角冰晶的折射形成，常见半径约为 $22^\circ$ 或 $46^\circ$；其圆心在太阳或月亮，而普通虹的圆心在反太阳点。
  \item “宝光”或 glory 出现在反太阳方向，涉及小液滴的后向散射、表面波和衍射，不能仅用一次或两次几何反射解释。
\end{itemize}
因此，几何光学能够给出许多弧带的主方向和色序，但小尺度粒子造成的展宽、干涉和衍射必须留到波动光学中处理。
\end{xesolution}

\Needspace{8\baselineskip}\begin{origexamplebox}{固定棱镜中的不同颜色}
\label{ch27:ex:depth-abbe}
光学玻璃的阿贝数定义为 $\nu=(n_D-1)/(n_F-n_C)$。顶角 $A=60^\circ$ 的棱镜有 $n_D=1.500$、$n_C=1.495$、$\nu=40$，置于空气中。白光以同一方向入射，调节棱镜，使 $D$ 光在棱镜内的传播方向垂直于顶角平分面。此后固定棱镜，求出射的 $C$ 光和 $F$ 光之间的夹角。
\end{origexamplebox}\par\nopagebreak[3]
\begin{xesolution}
先区分三种折射率。阿贝数给出的是两条谱线的折射率差，而不是折射率与波长的比值：
\[
 n_F-n_C=\frac{n_D-1}{\nu}=\frac{0.500}{40}=0.0125,
 \qquad n_F=1.5075.
\]
$D$ 光沿对称光路传播，因此在两个侧面处的内部折射角都是 $A/2=30^\circ$。参照图~\ref{ch27:fig:s42-19-6} 右图的两个工作面，分别相对该面法线计算折射角，第一面的共同入射角为
\[
 \sin i=n_D\sin(A/2)=0.750,
 \qquad i=48.5904^\circ.
\]
固定棱镜以后，$i$ 对三种颜色相同，但内部折射角并不相同。对任意折射率 $n$，逐次计算
\[
 r_1(n)=\arcsin\!\frac{\sin i}{n},\qquad
 r_2(n)=A-r_1(n),\qquad
 e(n)=\arcsin\!\bigl[n\sin r_2(n)\bigr].
\]
角 $e$ 从第二面外法线量起。两种颜色都从第二面透出，因为各自的 $n\sin r_2<1$；不能省略这一步存在性检查。代入得
\[
 e_C=48.1589^\circ,\qquad e_F=49.2438^\circ,
 \qquad\boxed{e_F-e_C=1.0849^\circ}.
\]
两束出射光共用同一个第二面法线，故出射角之差就是它们的夹角；总偏向角中的共同 $i-A$ 会抵消。$n_F>n_C$ 对应更大的偏向，与正常色散的方向一致。若令 $n_F=n_C$，两条计算链完全相同，夹角应为零。

本题 $A=60^\circ$，不能用薄棱镜式 $\delta\simeq(n-1)A$ 代替精确计算；也不能分别把 $C,F$ 两色都取为各自的最小偏向光路，因为实验中棱镜只调节了一次。
\end{xesolution}


\section{全反射、临界条件与光纤}



图~\ref{ch27:fig:s42-19-8} 的 A、B、C 依次为麦克斯韦、爱因斯坦和费曼：他们分别将光与电磁波、光量子和量子电动力学联系起来。D 为全反射实验情景，光从半圆形介质的圆弧面沿半径入射，再在平面边界反射；精确的角度条件以下面的临界角公式判定。
\historyfigure[0.88\linewidth]{optics-history-tir.png}{电磁理论与量子光学的代表人物，以及半圆形玻璃的全反射实验。\label{ch27:fig:s42-19-8}}

当光由高折射率介质 $n_1$ 射向低折射率介质 $n_2$，折射角随入射角增大得更快。折射角达到 $90^\circ$ 时的入射角称为临界角，满足
\begin{equation}
  \sin\theta_c=\frac{n_2}{n_1},\qquad n_1>n_2.
\end{equation}
当 $i>\theta_c$ 时不再存在传播的折射光，理想界面上发生全反射。

光纤利用芯层 $n_1$ 与包层 $n_2<n_1$ 的界面使光连续全反射，从而让传播方向跟随纤维弯曲。原理虽然简单，但长距离通信要求极低损耗。高锟指出早期光纤的主要损耗来自材料杂质，而非玻璃的不可避免吸收；高纯度材料由此把光纤变成实用通信媒介，他于 2009 年获得诺贝尔物理学奖。

全反射只限制芯层内光线在界面处的角度，并不意味着任意入射光都能被光纤接收。若外界折射率为 $n_0$，子午光线相对光纤轴的最大接收角 $\theta_a$ 满足
\[
 n_0\sin\theta_a=\sqrt{n_1^2-n_2^2}.
\]
右端称为数值孔径。弯曲过急会使外弯侧的界面入射角降到临界角以下，光便从包层泄漏；杂质吸收、散射和弯曲损耗则决定能够传输的实际距离。图~\ref{ch27:fig:s42-19-9} 展示光纤应用，图~\ref{ch27:fig:s42-19-10} 左侧则把芯层、包层和反射位置分开，说明“导光”不是光在单一玻璃棒中无条件反弹。

\historyfigure[0.92\linewidth]{fiber-pioneer-built-in-v1.png}{高锟、低损耗光纤与长距离光通信。\label{ch27:fig:s42-19-9}}

\subsection{从界面临界角推到光纤接收角}
设直光纤芯层折射率为 $n_1$、包层为 $n_2<n_1$，端面外介质为 $n_0$。只讨论穿过光纤轴的子午光线。光从端面进入后与轴成角 $\alpha$；由于端面法线沿轴，端面折射满足
\[
 n_0\sin\theta=n_1\sin\alpha.
\]
在圆柱侧壁处，法线垂直于轴，所以侧壁入射角为 $i=\pi/2-\alpha$。全反射要求
\[
 \sin i=\cos\alpha>\frac{n_2}{n_1}.
\]
将 $\sin^2\alpha=1-\cos^2\alpha$ 代回端面关系，得到
\[
 n_0^2\sin^2\theta<n_1^2-n_2^2.
\]
因此 $\sqrt{n_1^2-n_2^2}$ 称为相应子午光线模型的数值孔径。若其与 $n_0$ 的比值超过1，几何接收角的上限只能是 $90^\circ$，不能对大于1的数取反正弦。等号给临界极限，实际稳定导光还需留出余量；光纤弯曲、端面倾斜与三维斜光线会改变上述简单几何。

这也给出物理检验：$n_1\to n_2$ 时，接收范围缩为零；增大芯层和包层的折射率差，接收范围变大。全反射是界面上不存在传播型折射波，不是说低折射率一侧的电磁场严格为零；后面的隐失场正用来补充这一点。


\FiberFermatPanel



\Needspace{7\baselineskip}\begin{origexamplebox}{【全反射】}\label{ch27:ex:s42-19-10}
为什么从玻璃砖的一个侧面通常看不到由相邻侧面射入砖内的光线？玻璃折射率最大为多少时，才可能从侧面看见相邻侧面入射的光？
\end{origexamplebox}\par\nopagebreak[3]
\begin{xesolution}
设光从空气经一个面进入玻璃，玻璃内相对该面法线的折射角为 $r$。相邻面法线与原法线垂直，所以光到相邻面时入射角为 $90^\circ-r$。要从相邻面射出，必须满足
\[
 90^\circ-r\le\theta_c,\qquad \sin\theta_c=\frac1n.
\]
而空气侧入射角最大为 $90^\circ$，故玻璃内最大折射角满足 $\sin r_{\max}=1/n$，即 $r_{\max}=\theta_c$。存在可出射光线的条件是 $90^\circ-\theta_c\le\theta_c$，从而
\[
 \theta_c\ge45^\circ,\qquad \boxed{n\le\sqrt2}.
\]
这里的等号只是双掠射极限：$n=\sqrt2$ 要求入射和出射的空气侧光线都沿表面掠行，不能提供有限的贯穿通光量。因此实际非零通光量的条件为 $1<n<\sqrt2$，上确界为 $\sqrt2$。常见玻璃 $n\approx1.5>\sqrt2$，相邻面上发生全反射，不能沿该路径直接看到入射光束。
\end{xesolution}

\section{费马原理与非均匀介质中的光路}



逐个界面应用直线传播、反射和折射定律，原则上可以追踪任意光线；但复杂曲面会使法线和面形设计变得困难。费马原理提供了统一表述：单色光的真实光路使光程相对于邻近光路取驻值；在色散可忽略时，也可表述为传播时间取驻值。对本讲中的常规问题，它通常表现为极小值。

在色散可忽略的静止介质中，一段几何路程 $\dd s$ 的传播时间为 $\dd t=n\dd s/c$。定义光程
\begin{equation}
  \mathcal L=\int n\,\dd s,\qquad t=\frac{\mathcal L}{c}.
\end{equation}
真实光路使 $\delta\mathcal L=0$。这一表述可推导反射定律和折射定律，也解释为什么光线宁可在低折射率介质中走较长的几何距离。若同一点发出的不同光线经光学系统后严格会聚于同一点，则各条光路的总光程必须相同，否则会聚点处相位不能一致。

以平界面为例，界面两侧固定起点与终点到界面的垂直距离分别为 $a,b$，光线在界面上的交点横坐标为 $x$，两点的水平间隔为 $L$。光程
\[
 \mathcal L(x)=n_1\sqrt{a^2+x^2}
 +n_2\sqrt{b^2+(L-x)^2}.
\]
令 $\dd\mathcal L/\dd x=0$，得到
\[
 n_1\frac{x}{\sqrt{a^2+x^2}}
 =n_2\frac{L-x}{\sqrt{b^2+(L-x)^2}},
\]
两边的几何比值分别就是 $\sin i$ 与 $\sin r$，于是恢复斯涅尔定律。若两段都在同一介质而界面改为镜面，用镜像点把折线路径展开成直线，便得到 $i=r$。因此直线传播、反射、折射并不是彼此孤立的三条经验规则，而是同一个驻光程原理在不同约束下的结果。

例如，物点 $O=(-s_o,0)$、像点 $I=(s_i,0)$，界面顶点为原点，左侧折射率为 $1$、右侧为 $n$。界面点 $P=(x,y)$ 必须满足
\begin{equation}
 \sqrt{(x+s_o)^2+y^2}+n\sqrt{(x-s_i)^2+y^2}=s_o+n s_i.
\end{equation}
这是一类高次非球面。它说明即使“点到点严格成像”这一基本要求，也往往需要复杂且高精度的曲面；人的角膜与晶状体系统则是自然形成的精密光学系统。





\subsection{驻值方程如何成为折射定律}
将平界面取为横轴，起点在 $(0,a)$、终点在 $(L,-b)$，两侧折射率为 $n_1,n_2$，其中 $a,b>0$。经过界面点 $(x,0)$ 的候选折线路径，其光程为
\[
 \mathcal L(x)=n_1\sqrt{a^2+x^2}+n_2\sqrt{b^2+(L-x)^2}.
\]
求导时，第二段中的 $L-x$ 对 $x$ 求导为 $-1$，因此
\[
 \mathcal L'(x)=n_1\frac{x}{\sqrt{a^2+x^2}}
              -n_2\frac{L-x}{\sqrt{b^2+(L-x)^2}}.
\]
两个分式分别是两侧光线与法线夹角的正弦。令 $\mathcal L'=0$ 就得到斯涅尔定律，而不是一个独立的新折射公式。再求导，
\[
 \mathcal L''(x)=\frac{n_1a^2}{(a^2+x^2)^{3/2}}
                +\frac{n_2b^2}{[b^2+(L-x)^2]^{3/2}}>0,
\]
所以这一具体问题的驻值确实是唯一极小值。一般光学系统中的驻值不一定是最短路；不能把这里的正二阶导数结论推广到所有反射面、焦点和受约束光路。

分层介质中，每一薄层界面均平行，故相对层面法线的角 $i$ 满足 $n\sin i=\text{常量}$。如果改用相对层面的角 $\alpha$，同一规律写成 $n\cos\alpha=\text{常量}$。两式并不矛盾，差别只是角度的起算方向。建立微分方程前，应把图中的角与公式中的角核对一次。


% Relocated FisherBayDiagram to its own solution.

\Needspace{7\baselineskip}\begin{origexamplebox}{【运动学？】}\label{ch27:ex:s42-19-6}
湖湾由从顶点 $D$ 出发的两条直岸围成，夹角为 $\alpha$，$0<\alpha<\pi$。渔人的房子 $A$ 在一岸的陆地一侧，$C$ 是 $A$ 到该岸的垂足，$AC=h$，$AD=l$，且 $0<h<l$、$C$ 在岸线射线上。朋友的房子 $B$ 关于湖湾角平分线与 $A$ 对称。陆上步行速度为 $v$，划船速度为 $v/2$，可在岸上任意位置上下船。规划从 $A$ 到 $B$ 的最省时路线。
\end{origexamplebox}\par\nopagebreak[3]
\begin{xesolution}
\FisherBayDiagram

先分清陆地和水域。图~\ref{ch27:fig:s42-19-15} 中 $DC$ 是一侧岸线，$A,B$ 都不在水中；因此不能从房子直接划船。令
\[
 c=DC=\sqrt{l^2-h^2},\qquad \beta=\alpha/2.
\]
两岸上船、下船点记为 $P,Q$，$DP=t_1,DQ=t_2$。直线步行段均在各自陆地半平面内，船的直线段在湖湾内。其时间为
\[
 vT=\sqrt{h^2+(t_1-c)^2}+\sqrt{h^2+(t_2-c)^2}
 +2\sqrt{t_1^2+t_2^2-2t_1t_2\cos\alpha}.
\]
为什么可以取对称路线？前两项是严格凸函数，第三项是线性坐标差的欧氏范数，也是凸函数；交换 $t_1,t_2$ 不改变时间。将一条路线与它的对称路线取参数平均，凸性给出
$T((t_1+t_2)/2,(t_1+t_2)/2)\le T(t_1,t_2)$。
所以最优点必有 $t_1=t_2=t\ge0$。这不是仅凭图形对称猜测。
此时 $PQ=2t\sin\beta$，有
\[
 T(t)=\frac2v\left[\sqrt{h^2+(t-c)^2}+2t\sin\beta\right],
 \quad
 \frac v2T'(t)=\frac{t-c}{\sqrt{h^2+(t-c)^2}}+2\sin\beta.
\]
令 $k=2\sin\beta$，且
$T''(t)=2h^2/[v(h^2+(t-c)^2)^{3/2}]>0$。
由于 $T'(0)=2(k-c/l)/v$，答案完整分为
\[
 \boxed{\begin{cases}
 k\ge c/l:&t_*=0,\quad T_{\min}=2l/v;\\[2pt]
 0<k<c/l:&t_*=c-\dfrac{kh}{\sqrt{1-k^2}},\quad
 T_{\min}=\dfrac2v\left[h\sqrt{1-k^2}+kc\right].
 \end{cases}}
\]
第一种情形沿陆地 $A\to D\to B$ 行走，不划船；第二种情形沿 $A\to P$ 步行、$P\to Q$ 划船、$Q\to B$ 步行。
后一条件自动保证 $k<1$ 和 $t_*>0$。两式在 $k=c/l$ 处连续，时间单位均为长度除以速度；$h$ 不由 $l\sin\alpha$ 决定。
同岸的多次上下船不会更快，因为把这段船路改为沿同岸直线步行既不更长、速度又更大。于是上述一次横渡或完全绕岸的比较已覆盖最优路线。
它体现了费马最短时间思想：总距离最短不等于总时间最短。
\end{xesolution}

\Needspace{7\baselineskip}\begin{origexamplebox}{【严格聚焦】}\label{ch27:ex:s42-19-7}
球面镜不能把所有平行于轴的入射光严格会聚于一点。能够完成这一功能的反射面应具有何种形状？
\end{origexamplebox}\par\nopagebreak[3]
\begin{xesolution}
取顶点为原点、轴向右为 $x$ 正方向，焦点 $F=(f,0)$，$f>0$。平行光必须从右向左入射，照到向右开口的反射面。参考平面取 $x=a$，位于整个镜面右侧；反射点为 $P=(x,y)$。轴上光路长为 $a+f$，等光程条件为
\[
 (a-x)+\sqrt{(x-f)^2+y^2}=a+f.
\]
故 $PF=x+f$。平方并消去 $x^2+f^2$，得到 $y^2=4fx$；它也表明 $PF$ 等于 $P$ 到准线 $x=-f$ 的距离。因此子午截面是抛物线，旋转后为抛物面：
\[
 \boxed{y^2=4fx}
\]
也可直接检验反射定律。曲面法向可取 $\boldsymbol N=(-2f,y)$，入射单位方向为 $\boldsymbol k=(-1,0)$。反射方向
\[
 \boldsymbol k'=\boldsymbol k-2\frac{\boldsymbol k\cdot\boldsymbol N}{\boldsymbol N\cdot\boldsymbol N}\boldsymbol N
 =\frac{(4f^2-y^2,-4fy)}{4f^2+y^2}
 =\frac{F-P}{|F-P|}.
\]
所以反射光真正射向 $F$，而非仅延长线经过它。抛物面反射镜对轴上平行光严格无球差；有限口径带来衍射，制造误差和离轴入射仍会带来像差。
\end{xesolution}

\Needspace{7\baselineskip}\begin{origexamplebox}{【变折射率折射】}\label{ch27:ex:s42-19-5}
介质位于 $x>0,y>0$，折射率仅随 $y$ 连续变化。一束细光束沿 $x$ 方向垂直入射，并沿半径为 $R$ 的圆弧传播。求 $n(y)$。已知 $n(0)=n_0=1$，材料最大折射率为 $2.5$，求圆弧允许的最大圆心角。
\end{origexamplebox}\par\nopagebreak[3]
\begin{xesolution}
\GradientCircleDiagram

%<optics-review:gradient-geometry>
先把圆弧的角度与高度联系起来。以入射点为 $(0,0)$、圆心为 $(0,R)$，沿光的传播方向用参数 $\theta$ 写圆弧：
\[
 x=R\sin\theta,\qquad y=R(1-\cos\theta).
\]
微分得到 $(\dd x,\dd y)=R(\cos\theta,\sin\theta)\dd\theta$，所以 $\theta$ 同时就是切线与水平线的夹角。这便给出后面所用的 $\cos\theta=1-y/R$。水平薄层的法线竖直，斯涅尔定律中的入射角为 $\pi/2-\theta$，正弦因而变为 $\cos\theta$，不能在此误写成 $n\sin\theta$ 守恒。允许轨迹在 $0\le\theta<\pi/2$ 内；材料上限最终把高度限制到 $0\le y\le0.6R$，尚未到圆弧最高点。

%</optics-review:gradient-geometry>

把介质分成许多水平薄层。斯涅尔定律在各层界面上连续应用，得到守恒量 $n(y)\cos\theta=n_0$，其中 $\theta$ 是光线切线与 $x$ 轴的夹角。对图~\ref{ch27:fig:s42-19-14} 的圆弧，圆心在入射点正上方，有
\[
 \cos\theta=1-\frac{y}{R}.
\]
因此
\[
 \boxed{n(y)=\frac{n_0}{1-y/R}}.
\]
材料条件 $n\le2.5$ 给出 $\cos\theta\ge1/2.5=0.4$，故
\[
 \boxed{\theta_{\max}=\arccos0.4\approx66.4^\circ}.
\]
该角正是从入射半径转到当前半径的圆心角；若把它错当成相对法线的角，会得到互余的错误答案。
\end{xesolution}

\Needspace{7\baselineskip}\begin{origexamplebox}{【圆形光线】}\label{ch27:ex:s42-19-8}
一颗半径为 $R$ 的行星，其大气折射率随高度 $h$ 线性变化：$n=n_0-\alpha h$，其中 $n_0$ 是表面折射率。若某高度上沿行星表面切向发出的光恰能绕行星作圆周传播，求该高度 $h_0$。
\end{origexamplebox}\par\nopagebreak[3]
\begin{xesolution}
%<optics-review:circular-ray-proof>
关键是先解释圆轨道条件从何而来。设 $s$ 是沿光路的弧长，$\boldsymbol t=\dd\boldsymbol r/\dd s$ 为单位切向量。固定光路两端，只让中间各点有微小位移 $\delta\boldsymbol r$。光程 $\mathcal L=\int n\,\dd s$ 的一阶变化为
\[
 \delta\mathcal L=\int\!\left(\nabla n\cdot\delta\boldsymbol r
       +n\boldsymbol t\cdot\frac{\dd(\delta\boldsymbol r)}{\dd s}\right)\dd s.
\]
第二项分部积分，端点项因两端固定而为零。任意中间位移都要满足驻值条件，故
\[
 \frac{\dd(n\boldsymbol t)}{\dd s}=\nabla n.
\]
这说明折射率梯度怎样使光线弯曲。对半径为 $r$ 的圆，$n$ 沿圆周不变，单位切向量每走弧长 $\dd s$ 就向圆心转过 $\dd s/r$，即
\[
 \frac{\dd\boldsymbol t}{\dd s}=-\frac{\boldsymbol e_r}{r},
 \qquad \nabla n=\frac{\dd n}{\dd r}\boldsymbol e_r.
\]
因此 $-n/r=\dd n/\dd r$，整理即 $\dd(nr)/\dd r=0$。这个条件来自光线的局部弯曲规律；只说“沿圆周 $nr$ 为常量”还不足以推出它。

%</optics-review:circular-ray-proof>

球对称介质的射线守恒量为 $nr\sin\psi$，$\psi$ 是光线与径向的夹角。圆形轨道要求 $\psi=90^\circ$，并且 $n(r)r$ 在该半径处取驻值：
\[
 \frac{\dd}{\dd r}[n(r)r]=0,\qquad r=R+h.
\]
代入 $n=n_0-\alpha h$ 得
\[
 n_0-\alpha h-\alpha(R+h)=0,\qquad
 \boxed{h_0=\frac{n_0-\alpha R}{2\alpha}}.
\]
结果还必须满足 $h_0\ge0$ 且 $n(h_0)>0$；若材料参数不满足这些条件，物理解不存在。

%<optics-review:circular-ray-domain>
若 $n_0>0$，有限半径圆轨道要求折射率向外减小，即 $\alpha>0$；$\alpha=0$ 时光在均匀介质中沿直线传播，不能代入含 $1/\alpha$ 的结果。此时 $h_0\ge0$ 等价于 $n_0\ge\alpha R$，并且
\[
 r_0=\frac{n_0+\alpha R}{2\alpha},\qquad
 n(h_0)=\frac{n_0+\alpha R}{2}=\alpha r_0>0.
\]
还应确认算出的高度处于题设线性大气模型的适用范围；若要求严格在表面上方，则需 $n_0>\alpha R$。$n$ 无量纲，$\alpha$ 的量纲为长度的倒数，所得 $h_0$ 确为长度。本题求的是恰好切向发射的圆轨道，不意味着偏离该条件后仍能绕圆传播。

%</optics-review:circular-ray-domain>
\end{xesolution}

\section{界面上的电磁波与隐失场（选读）}

\subsection{反射、透射与偏振分量}


光线模型给出反射角和折射角，却不能回答反射光有多强、相位是否突变，以及为何反射率与偏振有关。
设折射率为 $n_1$ 的非磁性介质中的平面波，以入射角 $\theta_i$ 射向折射率为 $n_2$ 的介质，折射角为 $\theta_t$。
界面平移对称使频率和切向波矢连续，因而
\begin{equation}
 n_1\sin\theta_i=n_2\sin\theta_t.
\end{equation}
入射光、反射光和法线所在平面称入射面。电场垂直入射面的模式称 s 偏振，电场平行入射面的模式称 p 偏振。

\begin{center}\begin{minipage}{\linewidth}\centering
\begin{tikzpicture}[scale=.92,>=Latex]
  \fill[DeepBlue!6] (-4,-2.55) rectangle (4,0);
  \draw[very thick] (-4,0)--(4,0) node[right] {界面};
  \draw[dashed,guide] (0,-2.45)--(0,2.65) node[above] {法线};
  \coordinate (O) at (0,0);
  \draw[->,DeepBlue,very thick] (-3.0,2.2)--(O);
  \draw[->,curve,very thick] (O)--(3.0,2.2);
  \draw[->,DeepBlue,very thick] (O)--(1.55,-2.30);
  \node[above] at (-2.3,2.35) {入射};
  \node[above] at (2.4,2.35) {反射};
  \node[right] at (1.8,-2.3) {折射};
  \draw[->,guide] (0,.78) arc[start angle=90,end angle=144,radius=.78];
  \node at (-.48,.92) {$\theta_i$};
  \draw[->,guide] (0,-.80) arc[start angle=-90,end angle=-56,radius=.80];
  \node[left] at (-.65,-1.35) {$\theta_t$};
  \draw[guide,thin] (-.60,-1.32)--(.22,-.62);
  \node[left] at (-3.85,1.72) {$n_1$};
  \node[left] at (-3.85,-1.72) {$n_2$};
  \node[DeepBlue] at (-2.55,-.55) {$\odot\;s$：电场垂直纸面};
  \draw[->,DeepBlue,thick] (2.15,-1.05)--(2.78,-.62);
  \node[DeepBlue,right] at (2.12,-1.22) {$p$：电场在入射面内};
\end{tikzpicture}
\figcap{s、p 偏振的定义依赖入射面，而不是依赖纸张本身。图中折射光向法线偏折，对应 $n_2>n_1$。}
\end{minipage}\end{center}

规定界面法向 $\boldsymbol N$ 指向第二介质，$\boldsymbol s$ 垂直入射面，对入射、反射、透射各束分别取
$\boldsymbol p_j=\boldsymbol s\times\hat{\boldsymbol k}_j$。
这一右手构造使反射 p 基的切向方向与入射 p 基相反；因此振幅系数的正负必须连同基矢理解。
无自由面电荷、面电流且介质非磁性时，切向 $\boldsymbol E,\boldsymbol H$ 连续。用 $H=nE/Z_0$，s 分量满足
\[
 1+r_s=t_s,\qquad n_1\cos\theta_i(1-r_s)=n_2\cos\theta_t\,t_s.
\]
消去 $t_s$，可得下列第一行。p 分量则满足
\[
 \cos\theta_i(1-r_p)=\cos\theta_t\,t_p,\qquad
 n_1(1+r_p)=n_2t_p.
\]
第二式给 $t_p=n_1(1+r_p)/n_2$，代回第一式即得第二行。这就是菲涅尔公式（Fresnel equations）：
\begin{align}
 r_s&=\frac{n_1\cos\theta_i-n_2\cos\theta_t}
 {n_1\cos\theta_i+n_2\cos\theta_t},&
 t_s&=\frac{2n_1\cos\theta_i}
 {n_1\cos\theta_i+n_2\cos\theta_t},\\
 r_p&=\frac{n_2\cos\theta_i-n_1\cos\theta_t}
 {n_2\cos\theta_i+n_1\cos\theta_t},&
 t_p&=\frac{2n_1\cos\theta_i}
 {n_2\cos\theta_i+n_1\cos\theta_t}.
\end{align}
这里 $r,t$ 是按各自 s、p 正交基定义的电场复振幅比；p 基在反射前后的指向约定会影响 $r_p$ 的符号，但不影响可观测反射率。
对无吸收介质，能流反射率和透射率为
\begin{equation}
 R_{s,p}=|r_{s,p}|^2,
 \qquad
 T_{s,p}=\frac{n_2\cos\theta_t}{n_1\cos\theta_i}|t_{s,p}|^2,
 \qquad R_{s,p}+T_{s,p}=1.
\end{equation}
透射率中不能漏掉折射率和光束投影面积因子。正入射且 $n_2>n_1$ 时，反射电场相对入射电场发生 $\pi$ 相位突变，
这就是薄膜干涉中常说的“半波损失”；由高折射率侧向低折射率侧反射则没有这一突变。

对 p 偏振，当 $r_p=0$ 时反射光消失。结合折射定律可得
\begin{equation}
 \theta_i+\theta_t=\frac\pi2,
 \qquad
 \boxed{\tan\theta_B=\frac{n_2}{n_1}}.
\end{equation}
$\theta_B$ 称布儒斯特角（Brewster angle）。自然光在该角附近反射后，以 s 偏振为主；气体激光器中的布儒斯特窗正利用
这一性质降低 p 偏振的反射损耗。

\subsection{全反射中的隐失场}


当 $n_1>n_2$ 且 $\theta_i>\theta_c=\arcsin(n_2/n_1)$ 时，第二介质中不再有向法线方向传播的平面波，
但界面下方的电磁场并不为零。取界面为 $y=0$，第二介质位于 $y<0$，切向波数为
\[
 k_x=k_0n_1\sin\theta_i,
\]
其中 $k_0=\omega/c=2\pi/\lambda_0$，$\lambda_0$ 是真空波长，统一采用时间因子 $e^{-\ii\omega t}$。由 $k_x^2+k_y^2=n_2^2k_0^2$，第二介质中的法向波数成为纯虚数：
\begin{equation}
 k_y=-\ii\kappa,
 \qquad
 \kappa=k_0\sqrt{n_1^2\sin^2\theta_i-n_2^2}.
\end{equation}
于是透入第二介质的场具有
\begin{equation}
 E_t\propto e^{\ii k_xx}e^{\kappa y},
 \qquad
 \delta=\frac1\kappa
\end{equation}
的形式，$\delta$ 是振幅衰减到 $1/e$ 的穿透深度，强度相应衰减到 $e^{-2}$。选择 $k_y=-\ii\kappa$ 是为了使 $y\to-\infty$ 时场衰减，而不是无源增长。远离临界角且折射率差不小时，$\delta$ 常与波长同量级；当 $\theta_i\to\theta_c^+$ 时，$\kappa\to0$、$\delta\to\infty$，不可概括为始终只有一个波长。

\begin{center}\begin{minipage}{\linewidth}\centering
\begin{tikzpicture}[scale=.90,>=Latex]
  \fill[DeepBlue!6] (-4,-2.55) rectangle (4,0);
  \draw[very thick] (-4,0)--(4,0);
  \draw[dashed,guide] (0,-2.35)--(0,2.55) node[above] {法线};
  \draw[->,DeepBlue,very thick] (-3.0,2.05)--(0,0);
  \draw[->,curve,very thick] (0,0)--(3.0,2.05);
  \draw[DeepBlue,very thick] plot[smooth,domain=0:-2.25,samples=80,variable=\y]
    ({.9*exp(1.15*\y)},\y);
  \draw[->,DeepBlue,thick] (1.5,-1.45)--(3.4,-1.45);
  \node at (2.5,-1.8) {沿界面相位传播};
  \node[anchor=west] at (1.3,-.45) {$|E(y)|\propto e^{\kappa y}$};
  \node[anchor=west] at (1.3,-.90) {横向偏移表示振幅};
  \draw[<->,guide] (-1.15,0)--(-1.15,-.87) node[midway,left] {$\delta$};
  \node[left] at (-3.75,1.50) {$n_1$};
  \node[left] at (-3.75,-1.60) {$n_2<n_1$};
\end{tikzpicture}
\figcap{全反射与法向振幅包络：低折射率介质内的实线不是光线轨迹，其横向偏移表示振幅，法向不发生正弦振荡。\label{ch27:fig:evanescent}}
\end{minipage}\end{center}

图\ref{ch27:fig:evanescent}中 $y=-\delta$ 的振幅为界面处的 $1/e$。对于无损、无限厚第二介质，隐失场沿法向没有平均净能流，却可沿界面输运能量和相位。若在距离界面数个 $\delta$ 内放入另一块高折射率介质，
隐失场可耦合成传播波，出现受抑全反射；从量子图像看，它与隧穿具有相同的指数衰减结构。

\begin{starread}{有限光束在全反射中的微小位移}
平面波全反射系数的模为 $1$，但相位 $\phi_{s,p}$ 随入射角变化。真实窄光束是邻近切向波数 $k_x$ 的角谱叠加，
各分量反射后得到略有差别的相位；用驻相条件求反射光束中心，得到古斯--汉欣位移（Goos--H\"anchen shift）
\[
 \Delta_{\rm GH}=-\frac{\partial\phi_{s,p}}{\partial k_x}.
\]
它表现为反射点沿入射面方向移动，数量级通常接近波长。若进一步区分圆偏振分量，还会出现垂直入射面的
安贝尔--费多罗夫位移（Imbert--Fedorov shift）；它与光的自旋角动量和界面处的角动量守恒有关。
\end{starread}



\section{程函与光线方程（选读）}



在折射率缓慢变化的介质中，把单色标量场写成
\[
 \Psi(\vect r)=A(\vect r)e^{\ii k_0S(\vect r)}.
\]
这里 $S$ 的量纲是长度，$k_0S$ 才是无量纲相位。在无损、各向同性且 $n$ 缓变的标量近似中，亥姆霍兹方程为
$\nabla^2\Psi+k_0^2n^2\Psi=0$。代入上式并除去指数因子，逐项得到
\[
 \nabla^2A+2\ii k_0\nabla A\cdot\nabla S+\ii k_0A\nabla^2S
 +k_0^2A[n^2-|\nabla S|^2]=0.
\]
当振幅和折射率变化尺度远大于波长时，最高阶 $k_0^2$ 项先满足零，得到程函方程（eikonal equation）
\begin{equation}
 |\nabla S|=n(\vect r).
\end{equation}
等相位面法线给出光线方向。令弧长参数为 $s$、单位切向为
$\vect t=\nabla S/n$。沿光线求导为 $\dd/\dd s=\vect t\cdot\nabla$，再用梯度场无旋的恒等式
\[
 (\nabla S\cdot\nabla)\nabla S
 =\tfrac12\nabla|\nabla S|^2=\tfrac12\nabla n^2=n\nabla n,
\]
便有 $\dd(n\vect t)/\dd s=(\nabla S\cdot\nabla)\nabla S/n=\nabla n$，即光线方程
\begin{equation}
 \boxed{\frac{\dd}{\dd s}(n\vect t)=\nabla n}.
\end{equation}
它的切向分量是恒等式，法向分量说明光线向折射率增大的方向弯曲。若 $n=n(y)$，沿 $x$ 方向的平移对称性给出
\begin{equation}
 n(y)\sin\theta=\text{常量},
\end{equation}
其中 $\theta$ 是光线与层面法线的夹角；这就是分层介质中的连续折射定律。

同一光线方程也可由费马原理（Fermat principle）得到：固定端点之间的实际光路使光程取平稳值，
\begin{equation}
 \delta\int_A^B n(\vect r)\,\dd s=0.
\end{equation}
“平稳”不等同于“必为最短”，反射和某些多路径系统中可能取得极大值或鞍点；正确判据是一阶变分为零。

\Needspace{7\baselineskip}\begin{origexamplebox}{梯度折射率介质中的正弦光线}
平面分层介质的折射率在轴附近为
\[
 n(y)=n_0\left(1-\frac12g^2y^2\right),
 \qquad |gy|\ll1,
\]
其中 $g$ 的量纲为长度的倒数。一条近轴光线从 $x=0$ 处以 $y(0)=y_0$、$y'(0)=\alpha_0$ 入射。
求其轨迹，并说明这种介质为何能周期性会聚光束。

\begin{center}\begin{minipage}{\linewidth}\centering
\begin{tikzpicture}[x=.90cm,y=1.0cm,>=Latex]
  \fill[DeepBlue!6] (0,-2.0) rectangle (9,2.0);
  \draw[axis,->] (0,0)--(9.35,0) node[right] {$x$};
  \draw[axis,->] (0,-2.05)--(0,2.25) node[above] {$y$};
  \foreach \yy in {-1.5,-1,-.5,.5,1,1.5}{\draw[guide,dashed] (0,\yy)--(9,\yy);}
  \draw[curve,very thick,->,domain=0:8.8,samples=160]
    plot (\x,{1.22*cos(60*\x)+.20*sin(60*\x)});
  \draw[curve,thick,densely dashed,->,domain=0:8.8,samples=160]
    plot (\x,{.70*cos(60*\x)+.20*sin(60*\x)});
  \node[left] at (-.15,1.22) {$y_0$};
  \fill[DeepBlue] (1.5,.2) circle (1.5pt);
  \node[above] at (4.5,2.02) {$n$ 在轴上最大};
  \draw[<->,guide] (0,-2.35)--(6,-2.35) node[midway,below] {$2\pi/g$};
\end{tikzpicture}
\figcap{两条初始倾角相同、初始高度不同的光线在四分之一周期处相交；纵向尺度放大以辨认弯曲。空间周期为 $2\pi/g$。\label{ch27:fig:grin}}
\end{minipage}\end{center}
\end{origexamplebox}\par\nopagebreak[3]
\begin{xesolution}
近轴时 $s\simeq x$、$\vect t\simeq(1,y')$。光线方程的 $y$ 分量给出
\[
 \frac{\dd}{\dd x}\left(n\frac{\dd y}{\dd x}\right)=\frac{\dd n}{\dd y}.
\]
保留到 $y$ 和倾角的一阶量，左边可把 $n$ 取为 $n_0$，右边
$\dd n/\dd y=-n_0g^2y$，所以
\[
 y''+g^2y=0.
\]
结合初始条件，
\begin{equation}
 \boxed{y(x)=y_0\cos(gx)+\frac{\alpha_0}{g}\sin(gx)}.
\end{equation}
图\ref{ch27:fig:grin}中两条光线具有相同的 $\alpha_0$，在 $x=\pi/(2g)$ 处均有 $y=\alpha_0/g$，所以会聚点不必在轴上。光线以 $2\pi/g$ 为空间周期反复穿过轴附近。若入射的是平行于轴的光束，各条光线 $\alpha_0=0$ 而 $y_0$ 不同，则在
$x=(\pi/2+j\pi)/g$ 处同时到达 $y=0$，形成周期会聚；若来自同一轴上物点，则 $y_0=0$、$\alpha_0$ 不同，在 $x=j\pi/g$ 再会聚。一般任意一束光并不自动在同一点相交，必须说明其初始位置和倾角之间的关系。$g$ 的量纲为 $\mathrm{m^{-1}}$，因此 $gx$ 无量纲，轨迹也通过量纲检查。
若离轴过远使 $|gy|$ 不再远小于 $1$，二次折射率近似与正弦结论都会失效。
\end{xesolution}

\endgroup